\section{数学}

\subsection{公共头文件}
\noindent{\small \texttt{Math/head.hpp}}\par

\begin{lstlisting}
#pragma once
#include <bits/stdc++.h>
using namespace std;
using u32 = unsigned int;
using u64 = unsigned long long;
using ll = long long;
using i128 = __int128;
constexpr int MOD = 998'244'353;
constexpr int N = 1e6 + 50;
\end{lstlisting}

\subsection{模运算类}
\noindent{\small \texttt{Math/NT/modint.hpp}}\par

\begin{lstlisting}
#pragma once
#include "../head.hpp"

// ok
template <int P> 
struct ModInt {
    u32 val;
    constexpr ModInt(ll x = 0) : val((x %= P) < 0 ? x + P : x) {}
    // 注意i128等类型需要先转成ll
    using M = ModInt;

    // 希望编译期计算的话以下运算符全部补上constexpr
    M &operator+=(const M &p) {
        if ((val += p.val) >= P) {
            val -= P;
        }
        return *this;
    }
    M &operator-=(const M &p) {
        if ((val += P - p.val) >= P) {
            val -= P;
        }
        return *this;
    }
    M &operator*=(const M &p) {
        val = u64(val) * p.val % P;
        return *this;
    }
    M &operator/=(const M &p) { return (*this) *= p.inv(); }
    M operator-() const { return M(val ? P - val : u32(0)); }
    M operator+(const M &p) const { return M(*this) += p; }
    M operator-(const M &p) const { return M(*this) -= p; }
    M operator*(const M &p) const { return M(*this) *= p; }
    M operator/(const M &p) const { return M(*this) /= p; }
    bool operator==(const M &p) const { return val == p.val; }
    bool operator!=(const M &p) const { return val != p.val; }
    M pow(ll n) const {
        if (n < 0) return inv().pow(-n);
        M ret(1), mul(val);
        while (n > 0) {
            if (n & 1) ret *= mul;
            mul *= mul;
            n >>= 1;
        }
        return ret;
    }
    M inv() const {
        int a = val, b = P, u = 1, v = 0, t;
        while (b > 0) {
            t = a / b;
            a -= t * b;
            u -= t * v;
            swap(a, b);
            swap(u, v);
        }
        return M(u);
    }
    constexpr static M raw(u32 x) {
        M ret;
        ret.val = x;
        return ret;
    }
    constexpr static int mod() {
        return P;
    }
};
using mint = ModInt<MOD>;
\end{lstlisting}

\subsection{模运算与整数整除工具}
\noindent{\small \texttt{Math/NT/tools.hpp}}\par

\begin{lstlisting}
#pragma once
#include "../head.hpp"
#include <cstdlib>

// ok
inline ll norm(ll x, ll mod) { 
    if (0 <= x && x < mod) {
        return x;
    }
    return ((x %= mod) < 0) ? (x + mod) : x; 
}
inline ll mul(ll a, ll b, ll mod) { 
    a = norm(a, mod), b = norm(b, mod);
    return (i128)a * b % mod;
}
inline ll add(ll a, ll b, ll mod) {
    a = norm(a, mod), b = norm(b, mod);
    return ((a += b) >= mod) ? a - mod : a;
}
inline ll sub(ll a, ll b, ll mod) {
    a = norm(a, mod), b = norm(b, mod);
    return ((a -= b) < 0) ? a + mod : a;
}
inline ll qkp(ll a, ll b, ll mod) {
    a = norm(a, mod);
    ll ret = 1 % mod;
    for (; b; b >>= 1, a = mul(a, a, mod)) if (b & 1) ret = mul(ret, a, mod); 
    return ret;
}

// 要求质数模, O(log a)
inline ll inv(ll a, ll mod_prime) { 
    return qkp(a, mod_prime - 2, mod_prime); 
}

template <typename T>
T div_floor(T a, T b) {
    T q = a / b, r = a % b;
    if (r != 0 && (r < 0) != (b < 0)) --q;
    return q;
}
template <typename T>
T div_ceil(T a, T b) {
    T q = a / b, r = a % b;
    if (r != 0 && (r < 0) == (b < 0)) ++q;
    return q;
}
\end{lstlisting}

\subsection{扩展欧几里得与不定方程}
\noindent{\small \texttt{Math/NT/exgcd.hpp}}\par

\begin{lstlisting}
#pragma once
#include "../head.hpp"
// ok
inline i128 abs128(i128 x) {
    return x < 0 ? -x : x;
}
// O(log min(a, b))
inline array<i128, 3> exgcd(i128 a, i128 b) {
    if (b == 0) {
        if (a >= 0) return {1, 0, a};
        else return {-1, 0, -a};
    }
    auto [x, y, g] = exgcd(b, a % b);
    return {y, x - a / b * y, g};
}

inline i128 norm128(i128 x, i128 mod) {
    mod = abs128(mod);
    return ((x %= mod) < 0) ? x += mod : x;
}

// O(log min(a, b)), 要求a, b, c不超过ll范围
// solve ax + by = c, 其中x为被规范化后的值. 即0 <= x < b/g
// 步长为(b/g, -a/g)
inline tuple<bool, i128, i128> solveLDE(i128 a, i128 b, i128 c) {
    if (a == 0 && b == 0) {
        return {c == 0, 0, 0};
    }

    auto [x, y, g] = exgcd(a, b);
    if (c % g != 0) {
        return {false, 0, 0};
    }

    i128 k = c / g;
    if (b == 0) {
        return {true, x * k, 0};
    }

    x = norm128(x * k, b / g);
    y = (c - a * x) / b;
    return {true, x, y};
}

inline ll inv_exgcd(ll a, ll mod) {
    auto [x, y, g] = exgcd(a, mod);
    if (g != 1) return -1;
    return (ll)norm128(x, mod);
}
\end{lstlisting}

\subsection{扩展中国剩余定理}
\noindent{\small \texttt{Math/NT/excrt.hpp}}\par

\begin{lstlisting}
#pragma once
#include "../head.hpp"
#include "exgcd.hpp"
// ok

// 依赖norm128, exgcd
// 要求模数大于0, lcm 在i128范围内. 返回 { 最小非负解, lcm }, 无解返回{ -1, -1 }
// O(nlog(max(m)))
inline array<i128, 2> excrt(const vector<ll>& b, const vector<ll>& m) {
    int n = b.size();
    i128 nowm = m[0], nowb = norm128(b[0], nowm);

    for (int i = 1; i < n; i++) {
        ll A = nowm % m[i], B = m[i];
        i128 C = (i128)b[i] - nowb % B;

        auto [y1, _, g] = exgcd(A, B);
        if (C % g != 0) {
            return {-1, -1};
        }

        i128 M = B / g;
        i128 y = norm128(y1, M) * norm128(C / g, M) % M;

        nowb += y * nowm;
        nowm *= M;
    }
    return {nowb, nowm};
}
\end{lstlisting}

\subsection{试除分解与约数枚举}
\noindent{\small \texttt{Math/NT/factor.hpp}}\par

\begin{lstlisting}
#pragma once
#include "../head.hpp"

// ok
inline vector<pair<ll, int>> prime_factor(ll x) {
    vector<pair<ll, int>> ret;
    for (int i = 2; i <= x / i; i++) {
        int e = 0;
        while (x % i == 0) {
            x /= i;
            e++;
        }
        if (e > 0) {
            ret.emplace_back(i, e);
        }
    }
    if (x > 1) {
        ret.emplace_back(x, 1);
    }
    return ret;
}

inline vector<ll> factor(ll x) {
    vector<ll> head, tail;
    for (ll i = 1; i * i <= x; i++) {
        if (x % i == 0) {
            head.emplace_back(i);
            if (i * i < x) {
                tail.emplace_back(x / i);
            }
        }
    }
    head.insert(head.end(), tail.rbegin(), tail.rend());
    return head;
}
\end{lstlisting}

\subsection{Miller-Rabin 素性测试}
\noindent{\small \texttt{Math/NT/primetest.hpp}}\par

\begin{lstlisting}
#pragma once
#include "../head.hpp"
#include "tools.hpp"
// ok
// 依赖 mul, qkp
// O(log n)
inline bool primetest(ll n) {
    if (n < 2 || n % 6 % 4 != 1) {
        return (n | 1) == 3;
    }
    ll s = __builtin_ctzll(n - 1), d = n >> s;
    for (ll a : {2, 325, 9375, 28178, 450775, 9780504, 1795265022}) {
        ll p = qkp(a % n, d, n), i = s;
        while (p != 1 && p != n - 1 && a % n && i--) {
            p = mul(p, p, n);
        }
        if (p != n - 1 && i != s) {
            return false;
        }
    }
    return true;
}
\end{lstlisting}

\subsection{线性筛}
\noindent{\small \texttt{Math/NT/sieve.hpp}}\par

\begin{lstlisting}
#pragma once
#include "../head.hpp"
// ok
constexpr int PIN = 78'498;
// 1e8 : 5'761'455
// 1e7 : 664'579
// 1e6 : 78'498
// 1e5 : 9'592
// 1e4 : 1'229
// 1e3 : 168
// N>61000时 不超过N / (log(N) - 1.1)
inline bitset<N + 5> not_prime;
inline int primes[PIN + 5], pcnt = 0;
// O(N)
inline void sieve() {
    not_prime[0] = not_prime[1] = 1;
    for (ll i = 2; i <= N; i++) {
        if (!not_prime[i]) {
            primes[pcnt] = i;
            pcnt++;
        }
        for (int j = 0; j < pcnt && i * primes[j] <= N; j++) {
            not_prime[i * primes[j]] = 1;
            if (i % primes[j] == 0) {
                break;
            }
        }
    }
}
\end{lstlisting}

\subsection{扩展线性筛}
\noindent{\small \texttt{Math/NT/sieve\_ext.hpp}}\par

\begin{lstlisting}
#pragma once
#include "../head.hpp"
// ok
constexpr int PIN = 78'498;
inline bitset<N + 5> not_prime;
inline int primes[PIN + 5], pcnt = 0;
// 1e8 : 5'761'455
// 1e7 : 664'579
// 1e6 : 78'498
// 1e5 : 9'592
// 1e4 : 1'229
// 1e3 : 168
// N > 61000时不超过 N / (log(N) - 1.1)

inline int minp[N + 5]; // 最小质因子
// int phi[N + 5]; // 欧拉函数
// int mu[N + 5]; // 莫比乌斯函数
// int d[N + 5], times[N + 5]; // 约数计数函数
// ll sigma[N + 5]; int g[N + 5]; // 约数和函数

// O(N)
inline void sieve() {
    not_prime[0] = not_prime[1] = 1;
    // phi[1] = 1;
    // mu[1] = 1;
    // d[1] = 1;
    // sigma[1] = 1;
    for (ll i = 2; i <= N; i++) {
        if (!not_prime[i]) {
            primes[pcnt] = i, pcnt++;
            minp[i] = i;
            // phi[i] = i - 1;
            // mu[i] = -1;
            // d[i] = 2, times[i] = 1;
            // sigma[i] = i + 1, g[i] = i + 1;
        }
        for (int j = 0; j < pcnt && i * primes[j] <= N; j++) {
            int pj = primes[j];
            not_prime[i * pj] = 1;
            minp[i * pj] = pj;
            if (i % primes[j] == 0) {
                // phi[i * pj] = phi[i] * pj;
                // mu[i * pj] = 0;
                // d[i * pj] = d[i] / (times[i] + 1) * (times[i] + 2), times[i * pj] = times[i] + 1;
                // g[i * pj] = g[i] * pj + 1, sigma[i * pj] = sigma[i] / g[i] * g[i * pj];
                break;
            }
            // phi[i * pj] = phi[i] * (pj - 1);
            // mu[i * pj] = -mu[i];
            // d[i * pj] = d[i] * 2, times[i * pj] = 1;
            // sigma[i * pj] = sigma[i] * sigma[pj], g[i * pj] = pj + 1;
        }
    }
}

// O(log x), requires 1 <= x <= N
inline vector<pair<ll, int>> prime_factor(ll x) {
    vector<pair<ll, int>> ret;
    while (x > 1) {
        if (ret.empty() || ret.back().first != minp[x]) {
            ret.push_back({ minp[x], 1 });
        } else {
            ret.back().second++;
        }
        x /= minp[x];
    }
    return ret;
}
\end{lstlisting}

\subsection{分段筛}
\noindent{\small \texttt{Math/NT/sieve\_seg.hpp}}\par

\begin{lstlisting}
#pragma once
#include "../head.hpp"
// ok
constexpr ll R = 1'000'000'300;
constexpr int SQR = 32'000;
constexpr int PIS = 4'000; // pi(SQR)
// sqrt(10) < 3.17
inline bitset<SQR + 5> not_prime;
inline int primes[PIS + 5], pcnt = 0;
// O(SQRloglogSQR)
inline void init_sieve() {
    not_prime[0] = not_prime[1] = 1;
    for (ll i = 2; i * i <= R; i++) {
        if (!not_prime[i]) {
            primes[pcnt] = i;
            pcnt++;
            for (ll j = i * i; j <= SQR; j += i) {
                not_prime[j] = 1;
            }
        }
    }
}

// O(PIS + (r-l)loglogr)
inline vector<ll> get_primes(ll l, ll r) {
    if (l >= r) {
        return {};
    }
    r--;
    vector<bool> notp(r - l + 1);
    if (l == 0) {
        notp[0] = 1;
        if (r >= 1) {
            notp[1] = 1;
        }
    } else if (l == 1) {
        notp[0] = 1;
    }

    for (int i = 0; i < pcnt; i++) {
        ll p = primes[i];
        if (p * p > r) {
            break;
        }
        ll st = max(p * p, (l + p - 1) / p * p);
        for (ll j = st; j <= r; j += p) {
            notp[j - l] = 1;
        }
    }
    vector<ll> ret;
    for (ll i = l; i <= r; i++) {
        if (!notp[i - l]) {
            ret.emplace_back(i);
        }
    }
    return ret;
}
\end{lstlisting}

\subsection{类欧几里得求和}
\noindent{\small \texttt{Math/NT/euclid\_like.hpp}}\par

\begin{lstlisting}
#pragma once
#include "../head.hpp"

// ok
// O(log m)
inline i128 floor_sum(ll n, ll m, ll a, ll b) {
    if (n == 0) return 0;

    i128 ret = 0;
    if (a >= m) {
        ret += (i128)n * (n - 1) / 2 * (a / m);
        a %= m;
    }
    if (b >= m) {
        ret += (i128)(b / m) * n;
        b %= m;
    }
    if (a == 0) return ret;

    ll Y = ((i128)a * (n - 1) + b) / m;
    if (Y == 0) return ret;

    ret += (i128)n * Y - floor_sum(Y, a, m, m + a - b - 1);
    return ret;
}
\end{lstlisting}

\subsection{BSGS 与扩展 BSGS}
\noindent{\small \texttt{Math/NT/BSGS.hpp}}\par

\begin{lstlisting}
#pragma once
#include "../head.hpp"
#include "tools.hpp"
#include "exgcd.hpp"
#include "../../DataStructure/hashtable.hpp"

// ok

// 依赖norm, mul, qkp, exgcd, inv_exgcd, HashMap
// O(sqrt(m)), require gcd(a, m) = 1, 注意空间开销
inline ll bsgs(ll a, ll b, ll m) {
    a = norm(a, m);
    b = norm(b, m);

    if (m == 1) return 0;
    if (b == 1) return 0;

    ll n = sqrtl(m) + 1;
    HashMap<ll, ll> mp;

    // b * a^q
    ll cur = b;
    for (ll q = 0; q < n; q++) {
        mp[cur] = q;
        cur = mul(cur, a, m);
    }

    ll an = qkp(a, n, m);
    cur = 1;
    for (ll p = 1; p <= n + 1; p++) {
        cur = mul(cur, an, m);
        auto it = mp.find(cur);
        if (it != mp.end()) {
            ll x = p * n - it->second;
            if (x >= 0) return x;
        }
    }
    return -1;
}

// O(sqrt(m) + log(m))
inline ll exbsgs(ll a, ll b, ll m) {
    a = norm(a, m);
    b = norm(b, m);

    if (m == 1) return 0;
    if (b == 1) return 0;

    ll cnt = 0;
    ll cur = 1;
    while (true) {
        ll g = gcd(a, m);
        if (g == 1) break;
        if (b % g != 0) return -1;

        b /= g;
        m /= g;
        cur = mul(cur, a / g, m);
        cnt++;
        if (cur == b) return cnt;
    }

    ll iv = inv_exgcd(cur, m);
    ll rhs = mul(b, iv, m);
    ll t = bsgs(a, rhs, m);
    if (t == -1) return -1;
    return cnt + t;
}
\end{lstlisting}

\subsection{Pollard-Rho 分解（参考）}
\noindent{\small \texttt{Math/NT/\_pollard\_rho.hpp}}\par

\begin{lstlisting}
#pragma once
#include "../head.hpp"
#include "primetest.hpp"
#include "_divisors.hpp"

namespace rho {
inline mt19937_64 rng(chrono::steady_clock::now().time_since_epoch().count());

// 内部函数：传入大于一的合数 n，返回它的一个非平凡因子 d，满足 $1<d<n$。
// 使用 Brent 迭代和分批最大公约数计算，期望时间约为 $O(n^{1/4})$。
inline ll rho(ll n) {
    if (n % 2 == 0) return 2;
    if (n % 3 == 0) return 3;

    constexpr int B = 128;

    while (true) {
        ll c = uniform_int_distribution<ll>(1, n - 1)(rng);
        ll y = uniform_int_distribution<ll>(0, n - 1)(rng);

        // 传入当前迭代值 x，返回 $(x^2+c)\bmod n$，使用减法避免加法溢出。
        auto f = [&](ll x) -> ll {
            ll t = mul(x, x, n);
            return t >= n - c ? t - (n - c) : t + c;
        };

        ll x = 0, ys = 0;
        ll d = 1, r = 1;

        while (d == 1) {
            x = y;
            for (ll i = 0; i < r; i++) {
                y = f(y);
            }

            for (ll k = 0; k < r && d == 1; k += B) {
                ys = y;
                ll q = 1;
                ll len = min<ll>(B, r - k);

                for (ll i = 0; i < len; i++) {
                    y = f(y);
                    q = mul(q, abs(x - y), n);
                }

                d = gcd(q, n);
            }

            r *= 2;
        }

        if (d == n) {
            do {
                ys = f(ys);
                d = gcd(abs(x - ys), n);
            } while (d == 1);
        }

        if (d != n) return d;
    }
}

// 传入正整数 x，返回按质数递增排列的质因数分解，每项为 {质数, 指数}。
// x 须在有符号 ll 范围内；一返回空数组。外部调用名为 rho::prime_factor。
inline vector<pair<ll, int>> prime_factor(ll x) {
    assert(x >= 1);
    vector<ll> fac;

    // 递归分解传入的正整数 n，把所有质因子追加到 fac；self 用于递归，不直接返回结果。
    auto dfs = [&](auto self, ll n) -> void {
        if (n == 1) return;
        if (primetest(n)) {
            fac.push_back(n);
            return;
        }

        ll d = rho(n);
        self(self, d);
        self(self, n / d);
    };

    dfs(dfs, x);
    sort(fac.begin(), fac.end());

    vector<pair<ll, int>> ret;
    for (ll p : fac) {
        if (ret.empty() || ret.back().first != p) {
            ret.emplace_back(p, 1);
        } else {
            ret.back().second++;
        }
    }
    return ret;
}

// 传入正整数 x，返回递增排列的所有正约数；先调用 Pollard-Rho 分解，一返回 {1}。
inline vector<ll> factor(ll x) {
    return divisors(prime_factor(x));
}
} 
\end{lstlisting}

\subsection{由质因数分解生成约数（参考）}
\noindent{\small \texttt{Math/NT/\_divisors.hpp}}\par

\begin{lstlisting}
#pragma once
#include "../head.hpp"

// 传入正整数的质因数分解 fac，每项为 {质数, 指数}，返回递增排列的所有正约数。
// 空分解表示一，返回 {1}；原数须在 ll 范围内，时间为 $O(\tau(n)\log\tau(n))$。
inline vector<ll> divisors(const vector<pair<ll, int>>& fac) {
    vector<ll> ret{1};
    for (auto [p, e] : fac) {
        int n = ret.size();
        ll mul = 1;
        for (int k = 1; k <= e; k++) {
            mul *= p;
            for (int i = 0; i < n; i++) ret.push_back(ret[i] * mul);
        }
    }
    sort(ret.begin(), ret.end());
    return ret;
}
\end{lstlisting}

\subsection{莫比乌斯函数与欧拉函数筛（参考）}
\noindent{\small \texttt{Math/NT/\_mobius.hpp}}\par

\begin{lstlisting}
#pragma once
#include "../head.hpp"

// 线性筛出质数、最小质因子、欧拉函数和莫比乌斯函数，并保存后两者的前缀和。
// 构造时指定上限 n；结果依次存于 primes、minp、phi、mu、sphi、smu，下标从零到 n。
struct MobiusSieve {
    int n;
    vector<int> primes, minp, phi, mu;
    vector<ll> sphi, smu;
    // 传入正上限 n，建立筛表；一的欧拉函数和莫比乌斯函数均为一，零号项保留为零。
    // 预处理时间、空间均为 $O(n)$。
    MobiusSieve(int n) : n(n), minp(n + 1), phi(n + 1), mu(n + 1),
                         sphi(n + 1), smu(n + 1) {
        assert(n >= 1);
        phi[1] = mu[1] = 1;
        for (int i = 2; i <= n; i++) {
            if (!minp[i]) {
                primes.push_back(i);
                minp[i] = i;
                phi[i] = i - 1;
                mu[i] = -1;
            }
            for (int p : primes) {
                if (p > n / i) break;
                minp[i * p] = p;
                if (i % p == 0) {
                    phi[i * p] = phi[i] * p;
                    break;
                }
                phi[i * p] = phi[i] * (p - 1);
                mu[i * p] = -mu[i];
            }
        }
        for (int i = 1; i <= n; i++) {
            sphi[i] = sphi[i - 1] + phi[i];
            smu[i] = smu[i - 1] + mu[i];
        }
    }
};
\end{lstlisting}

\subsection{杜教筛（参考）}
\noindent{\small \texttt{Math/NT/\_dujiao\_sieve.hpp}}\par

\begin{lstlisting}
#pragma once
#include "_mobius.hpp"

// 预筛小范围并缓存大查询，用于求莫比乌斯函数、欧拉函数的前缀和。
// 构造时传入预筛上限，随后调用 getsmu 或 getsphi；预筛数组沿用 MobiusSieve。
struct DuJiao : MobiusSieve {
    unordered_map<ll, ll> mp_mu;
    unordered_map<ll, i128> mp_phi;
    // 传入正预筛上限 n，建立筛表和空查询缓存；预处理时间、空间均为 $O(n)$。
    DuJiao(int n = 1'000'000) : MobiusSieve(n) {}

    // 传入非负上限 x，返回从一到 x 的莫比乌斯函数之和，结果类型为 ll。
    ll getsmu(ll x) {
        assert(x >= 0);
        if (x <= n) return smu[x];
        auto it = mp_mu.find(x);
        if (it != mp_mu.end()) return it->second;
        i128 ret = 1;
        for (ll l = 2, r; l <= x; l = r + 1) {
            ll q = x / l;
            r = x / q;
            ret -= (i128)(r - l + 1) * getsmu(q);
            if (r == x) break;
        }
        return mp_mu[x] = (ll)ret;
    }
    // 传入非负上限 x，返回从一到 x 的欧拉函数之和，结果类型为 i128。
    i128 getsphi(ll x) {
        assert(x >= 0);
        if (x <= n) return sphi[x];
        auto it = mp_phi.find(x);
        if (it != mp_phi.end()) return it->second;
        i128 ret = (i128)x * ((i128)x + 1) / 2;
        for (ll l = 2, r; l <= x; l = r + 1) {
            ll q = x / l;
            r = x / q;
            ret -= (i128)(r - l + 1) * getsphi(q);
            if (r == x) break;
        }
        return mp_phi[x] = ret;
    }
};
\end{lstlisting}

\subsection{数论分块与带权求和（参考）}
\noindent{\small \texttt{Math/NT/\_quotient\_blocks.hpp}}\par

\begin{lstlisting}
#pragma once
#include "../head.hpp"

// 传入非负整数 n 和回调 work，按整数商分块依次调用 work(l, r, q)。
// 回调收到闭区间 [l, r] 及其中相同的商 q；本函数无返回值，块数为 $O(\sqrt n)$。
template <typename F>
void quotient_blocks(ll n, F work) {
    assert(n >= 0);
    for (ll l = 1, r; l <= n; l = r + 1) {
        ll q = n / l;
        r = n / q;
        work(l, r, q);
        if (r == n) break;
    }
}

// 传入非负上限 n、m，以及权函数 f 的前缀和 s，返回带权的整数商乘积之和：
// $\sum_{d=1}^{\min(n,m)}f(d)\lfloor n/d\rfloor\lfloor m/d\rfloor$，要求 s[0] 为零且长度大于 $\min(n,m)$。
// 结果类型与 s 的元素类型相同；较大整数结果用 i128，模意义结果用 mint。
template <typename Z>
Z quotient_sum(ll n, ll m, const vector<Z>& s) {
    assert(n >= 0 && m >= 0 && (ll)s.size() > min(n, m));
    Z ret = 0;
    ll lim = min(n, m);
    for (ll l = 1, r; l <= lim; l = r + 1) {
        ll a = n / l, b = m / l;
        r = min(n / a, m / b);
        ret += (s[r] - s[l - 1]) * Z(a) * Z(b);
        if (r == lim) break;
    }
    return ret;
}
\end{lstlisting}

\subsection{整除变换与最大公约数、最小公倍数卷积（参考）}
\noindent{\small \texttt{Math/NT/\_divisor\_transform.hpp}}\par

\begin{lstlisting}
#pragma once
#include "../head.hpp"

// 传入按正整数下标存放的数组 a，原地把每项替换为其所有约数位置的值之和。
// invert 为真时做逆变换；a[0] 不参与，不返回新数组，时间为 $O(n\log n)$。
template <typename Z>
void divisor_zeta(vector<Z>& a, bool invert = false) {
    int n = (int)a.size() - 1;
    if (!invert) {
        for (int i = n; i >= 1; i--) {
            for (int j = i + i; j <= n; j += i) a[j] += a[i];
        }
    } else {
        for (int i = 1; i <= n; i++) {
            for (int j = i + i; j <= n; j += i) a[j] -= a[i];
        }
    }
}

// 传入按正整数下标存放的数组 a，原地把每项替换为其所有倍数位置的值之和。
// invert 为真时做逆变换；a[0] 不参与，不返回新数组，时间为 $O(n\log n)$。
template <typename Z>
void multiple_zeta(vector<Z>& a, bool invert = false) {
    int n = (int)a.size() - 1;
    if (!invert) {
        for (int i = 1; i <= n; i++) {
            for (int j = i + i; j <= n; j += i) a[i] += a[j];
        }
    } else {
        for (int i = n; i >= 1; i--) {
            for (int j = i + i; j <= n; j += i) a[i] -= a[j];
        }
    }
}

// 传入等长、非空的数组 a、b，返回按最大公约数归类的卷积。
// 返回项为 $c_k=\sum_{\gcd(i,j)=k}a_i b_j$；输入下标从一开始，返回的零号项为零。
template <typename Z>
vector<Z> convolution_gcd(vector<Z> a, vector<Z> b) {
    assert(!a.empty() && a.size() == b.size());
    multiple_zeta(a);
    multiple_zeta(b);
    for (int i = 1; i < (int)a.size(); i++) a[i] *= b[i];
    multiple_zeta(a, true);
    a[0] = 0;
    return a;
}

// 传入等长、非空的数组 a、b，返回按最小公倍数归类的卷积。
// 返回项为 $c_k=\sum_{\operatorname{lcm}(i,j)=k}a_i b_j$，仅保留原数组范围内的下标；零号项为零。
template <typename Z>
vector<Z> convolution_lcm(vector<Z> a, vector<Z> b) {
    assert(!a.empty() && a.size() == b.size());
    divisor_zeta(a);
    divisor_zeta(b);
    for (int i = 1; i < (int)a.size(); i++) a[i] *= b[i];
    divisor_zeta(a, true);
    a[0] = 0;
    return a;
}
\end{lstlisting}

\subsection{Tonelli-Shanks 模平方根（参考）}
\noindent{\small \texttt{Math/NT/\_mod\_sqrt.hpp}}\par

\begin{lstlisting}
#pragma once
#include "../head.hpp"
#include "tools.hpp"

// 传入整数 a 和质数 p，求模平方根，返回较小的非负根；无解返回 -1。
// p 至少为二且须在有符号 ll 范围内；Tonelli-Shanks 主循环为 $O(\log^2 p)$，另需寻找二次非剩余。
inline ll mod_sqrt(ll a, ll p) {
    assert(p >= 2);
    a = norm(a, p);
    if (a == 0 || p == 2) return a;
    if (qkp(a, (p - 1) / 2, p) != 1) return -1;
    if (p % 4 == 3) {
        ll x = qkp(a, p / 4 + 1, p);
        return min(x, p - x);
    }
    ll d = p - 1;
    int s = __builtin_ctzll(d);
    d >>= s;
    ll z = 2;
    while (qkp(z, (p - 1) / 2, p) == 1) z++;
    ll x = qkp(a, (d + 1) / 2, p), t = qkp(a, d, p), c = qkp(z, d, p);
    while (t != 1) {
        int i = 0;
        ll cur = t;
        while (cur != 1) {
            cur = mul(cur, cur, p);
            i++;
        }
        ll b = c;
        for (int j = 0; j < s - i - 1; j++) b = mul(b, b, p);
        x = mul(x, b, p);
        c = mul(b, b, p);
        t = mul(t, c, p);
        s = i;
    }
    return min(x, p - x);
}
\end{lstlisting}

\subsection{exLucas 合数模组合数（参考）}
\noindent{\small \texttt{Math/NT/\_exlucas.hpp}}\par

\begin{lstlisting}
#pragma once
#include "../head.hpp"
#include "factor.hpp"
#include "tools.hpp"
#include "exgcd.hpp"

// 预处理合数模组合数；构造时传入正 int 模数，使用 nCr 查询，结果返回为 ll。
// 预处理时间、空间为 $O(\sum p^e)$，只适合各素数幂不大的模数。
struct ExLucas {
    // 保存一个素数幂模数及其去除素因子的阶乘表，供 ExLucas 内部使用。
    // 构造时指定质数 p 和正指数 e，查询结果都按该素数幂取模。
    struct Part {
        ll p, mod;
        int e;
        vector<ll> fact;
        // 传入质数 p 和正指数 e，预处理模 $p^e$ 的阶乘表，要求该素数幂不超过外层模数。
        Part(ll p, int e) : p(p), mod(1), e(e) {
            for (int i = 0; i < e; i++) mod *= p;
            fact.resize(mod + 1);
            fact[0] = 1;
            for (ll i = 1; i <= mod; i++) {
                fact[i] = fact[i - 1];
                if (i % p) fact[i] = mul(fact[i], i, mod);
            }
        }
        // 传入非负整数 n，返回 n 的阶乘中素因子 p 的指数。
        ll count(ll n) const {
            ll ret = 0;
            for (; n; n /= p) ret += n / p;
            return ret;
        }
        
        // 传入非负整数 n，返回从 n 的阶乘中去掉所有 p 因子后，剩余部分模当前 mod 的值。
        // 一个整块中所有与 p 互质的数之积只能是 $1$ 或 $-1$，因此只需检查整块数的奇偶。
        ll unit_fact(ll n) const {
            ll ret = 1;
            for (; n; n /= p) {
                if (n / mod & 1) ret = mul(ret, fact[mod], mod);
                ret = mul(ret, fact[n % mod], mod);
            }
            return ret;
        }
        // 内部查询：传入合法的 n、r，返回组合数模当前素数幂的值，要求 $0\le r\le n$。
        ll nCr(ll n, ll r) const {
            ll k = count(n) - count(r) - count(n - r);
            if (k >= e) return 0;
            ll denom = mul(unit_fact(r), unit_fact(n - r), mod);
            ll ret = mul(unit_fact(n), inv_exgcd(denom, mod), mod);
            return mul(ret, qkp(p, k, mod), mod);
        }
    };
    int mod;
    vector<Part> parts;
    // 传入正 int 模数 mod，分解模数并建立各素数幂的表；模数为一时无需建立表。
    ExLucas(int mod) : mod(mod) {
        assert(mod > 0);
        for (auto [p, e] : prime_factor(mod)) parts.emplace_back(p, e);
    }
    // 传入非负整数 n 和选取数 r，返回组合数模构造时指定的模数的最小非负值。
    // 非法 r 返回零；每个素数幂的查询时间约为 $O(\log_p n+\log M)$，M 是该素数幂模数。
    ll nCr(ll n, ll r) const {
        if (!(0 <= r && r <= n) || mod == 1) return 0;
        ll ret = 0;
        for (const auto& part : parts) {
            ll m = mod / part.mod;
            ret = (ret + (i128)part.nCr(n, r) * m * inv_exgcd(m, part.mod)) % mod;
        }
        return ret;
    }
};
\end{lstlisting}

\subsection{阶乘、组合数与 Lucas}
\noindent{\small \texttt{Math/Comb/factorial.hpp}}\par

\begin{lstlisting}
#pragma once
#include "../head.hpp"
#include "../NT/modint.hpp"

// ok
// 依赖 modint
// 要求N < MOD
template <int N, int MOD = 998'244'353>
struct Factorial {
    using mint = ModInt<MOD>;
    array<mint, N + 5> fact, factinv;
    Factorial() {
        fact[0] = 1;
        for (int i = 1; i <= N; i++) {
            fact[i] = fact[i - 1] * i;
        }
        factinv[N] = fact[N].inv();
        for (int i = N - 1; i >= 0; i--) {
            factinv[i] = factinv[i + 1] * (i + 1);
        }
    }
    mint nCr(ll n, ll r) {
        if (!(0 <= r && r <= n)) {
            return 0;
        }
        return fact[n] * factinv[r] * factinv[n - r];
    }
    mint nPr(ll n, ll r) {
        if (!(0 <= r && r <= n)) {
            return 0;
        }
        return fact[n] * factinv[n - r];
    }
};

// inline Factorial<N, MOD> f;

// O(min(r, n - r)), 要求r < MOD
template <int MOD = 998'244'353>
inline ModInt<MOD> nCr_large(ll n, ll r) {
    if (!(0 <= r && r <= n)) {
        return 0;
    }
    r = min(r, n - r);

    ModInt<MOD> ret = 1, denom = 1;
    for (ll i = 1; i <= r; i++) {
        ret *= n - i + 1;
        denom *= i;
    }
    return ret / denom;
};

// O(log_MOD n * C(n, r)) 其中C(n, r)为单次求组合数复杂度
template <int MOD = 998'244'353>
inline ModInt<MOD> lucas(ll n, ll r) {
    if (!(0 <= r && r <= n)) return 0;
    if (r == 0) {
        return 1;
    }

    return lucas<MOD>(n / MOD, r / MOD) * 
    nCr_large<MOD>(n % MOD, r % MOD); // 或改用其他求组合数方式
}
\end{lstlisting}

\subsection{斯特林数}
\noindent{\small \texttt{Math/Comb/stirling.hpp}}\par

\begin{lstlisting}
#pragma once
#include "../head.hpp"
#include "../NT/modint.hpp"
#include "../Poly/NTT.hpp"

// ok

// 依赖modint
template <int N, int MOD = 998'244'353>
struct Stirling {
    using mint = ModInt<MOD>;
    array<array<mint, N + 5>, N + 5> _c, _S;

    Stirling() {
        _c[0][0] = 1;
        _S[0][0] = 1;
        for (int n = 1; n <= N; n++) {
            for (int k = 1; k <= n; k++) {
                // c(n,k) = c(n-1,k-1) + (n-1)c(n-1,k)
                _c[n][k] = _c[n - 1][k - 1] + _c[n - 1][k] * (n - 1);

                // s(n,k) = s(n-1,k-1) - (n-1)s(n-1,k)

                // S(n,k) = S(n-1,k-1) + kS(n-1,k)
                _S[n][k] = _S[n - 1][k - 1] + _S[n - 1][k] * k;
            }
        }
    }
    
    mint c(ll n, ll k) {
        // 无符号第一类斯特林数
        if (k < 0 || k > n) {
            return 0;
        }
        return _c[n][k];
    }
    mint s(ll n, ll k) {
        // 有符号第一类斯特林数
        if (k < 0 || k > n) return 0;
        return (n - k) % 2 ? -_c[n][k] : _c[n][k];
    }
    mint S(ll n, ll k) {
        // 第二类斯特林数
        if (k < 0 || k > n) {
            return 0;
        }
        return _S[n][k];
    }
    
};

// inline Stirling<1'005> stir;



// 依赖NTT
// O(nlog n) 卷积求S(n), 要求n >= 0, 模数为NTT友好质数, need <= MOD
template <int MOD = 998'244'353>
vector<ModInt<MOD>> stirling2_row(ll n, ll need = -1) {
    // S(n, k) = 1/k! \sum_{i = 0}^k (-1)^{k - i} \binom{k}{i} i^n
    // 即(-1)^j / j! 与 i^n / i! 卷积
    using mint = ModInt<MOD>;

    if (need == -1 || need > n) {
        need = n + 1;
    }
    if (need == 0) return {};

    vector<mint> A(need), B(need);

    // A[i] = 1 / i!
    mint fac = 1;
    for (ll i = 1; i < need; i++) {
        fac *= i;
    }
    A[need - 1] = fac.inv();
    for (ll i = need - 1; i >= 1; i--) {
        A[i - 1] = A[i] * i;
    }

    // A[i] = (-1)^i / i!
    // B[i] = i^n / i!
    for (ll i = 0; i < need; i++) {
        B[i] = mint(i).pow(n) * A[i];
        if (i & 1) A[i] = -A[i];
    }

    return convolution(A, B, need);
}
\end{lstlisting}

\subsection{错排、拆分、Bell、Catalan 与树计数（参考）}
\noindent{\small \texttt{Math/Comb/\_numbers.hpp}}\par

\begin{lstlisting}
#pragma once
#include "../head.hpp"
#include "factorial.hpp"

// 传入非负上限 n，返回从零到 n 的错排数，按元素类型 Z 运算。
// 空排列计为一个错排；时间和空间均为 $O(n)$。
template <typename Z = mint>
vector<Z> derangements(int n) {
    assert(n >= 0);
    vector<Z> ret(n + 1);
    ret[0] = 1;
    for (int i = 2; i <= n; i++) ret[i] = (ret[i - 1] + ret[i - 2]) * Z(i - 1);
    return ret;
}

// 传入非负上限 n，返回从零到 n 的无序正整数拆分数，零的拆分数为一。
// 使用五边形数递推，时间为 $O(n\sqrt n)$，空间为 $O(n)$。
template <typename Z = mint>
vector<Z> partition_numbers(int n) {
    assert(n >= 0);
    vector<Z> ret(n + 1);
    ret[0] = 1;
    for (int i = 1; i <= n; i++) {
        for (ll j = 1; j * (3 * j - 1) / 2 <= i; j++) {
            ll a = j * (3 * j - 1) / 2, b = a + j;
            Z val = ret[i - a];
            if (b <= i) val += ret[i - b];
            ret[i] += j & 1 ? val : -val;
        }
    }
    return ret;
}

// 传入非负上限 n，返回从零到 n 的 Bell 数，即不同元素划分为无标号非空集合的方案数。
// 空集有一种划分；时间为 $O(n^2)$，额外空间为 $O(n)$。
template <typename Z = mint>
vector<Z> bell_numbers(int n) {
    assert(n >= 0);
    vector<Z> row(n + 1), ret(n + 1);
    row[0] = ret[0] = 1;
    for (int i = 1; i <= n; i++) {
        Z prev = row[0];
        row[0] = row[i - 1];
        for (int j = 1; j <= i; j++) {
            Z cur = row[j];
            row[j] = row[j - 1] + prev;
            prev = cur;
        }
        ret[i] = row[0];
    }
    return ret;
}

// 传入非负下标 n 和现有阶乘表 f，返回第 n 个 Catalan 数，要求 $2n\le N<P$。
// P 须为质数，返回值类型为 ModInt<P>；零号 Catalan 数为一。
template <int N, int P>
ModInt<P> catalan(ll n, Factorial<N, P>& f) {
    assert(0 <= n && n <= N / 2);
    return f.nCr(2 * n, n) - f.nCr(2 * n, n - 1);
}

// 传入正步数 a、负步数 b 和阶乘表 f，返回所有前缀和非负的步序列数量。
// a、b 均非负，要求 $a+b\le N<P$ 且 P 为质数；当 $a<b$ 时返回零。
template <int N, int P>
ModInt<P> ballot(ll a, ll b, Factorial<N, P>& f) {
    assert(a >= 0 && b >= 0 && a <= N && b <= N - a);
    if (a < b) return 0;
    return f.nCr(a + b, b) - f.nCr(a + b, b - 1);
}
// 传入各类元素数量 cnt 和阶乘表 f，返回多重集排列数。
// 数量均非负，要求总数不超过 N 且 $N<P$，P 为质数；空多重集返回一。
template <int N, int P>
ModInt<P> multinomial(const vector<int>& cnt, Factorial<N, P>& f) {
    int n = 0;
    ModInt<P> ret = 1;
    for (int x : cnt) {
        assert(x >= 0 && x <= N - n);
        n += x;
        ret *= f.factinv[x];
    }
    return ret * f.fact[n];
}

// 传入正点数 n，返回 n 个标号点的无根树数量，按类型 Z 运算；单点树返回一。
template <typename Z = mint>
Z cayley(ll n) {
    assert(n >= 1);
    return n == 1 ? Z(1) : Z(n).pow(n - 2);
}

// 传入各点度数 deg 和阶乘表 f，返回满足这些度数的标号树数量；无此类树时返回零。
// 至少有一个点；n 为点数，要求 $n-2\le N<P$ 且 P 为质数，单点树的合法度数为零。
template <int N, int P>
ModInt<P> tree_degree_count(const vector<int>& deg, Factorial<N, P>& f) {
    int n = deg.size();
    assert(n >= 1 && n - 2 <= N);
    if (n == 1) return deg[0] == 0;
    ll sum = 0;
    ModInt<P> ret = f.fact[n - 2];
    for (int d : deg) {
        if (d <= 0 || d >= n) return 0;
        sum += d;
        ret *= f.factinv[d - 1];
    }
    return sum == 2LL * n - 2 ? ret : ModInt<P>(0);
}
\end{lstlisting}

\subsection{Burnside、项链与手链（参考）}
\noindent{\small \texttt{Math/Comb/\_burnside.hpp}}\par

\begin{lstlisting}
#pragma once
#include "../head.hpp"
#include "../NT/modint.hpp"
#include "../NT/factor.hpp"
#include "../NT/_divisors.hpp"

// 传入完整置换群 group 和颜色数 colors，返回把群作用视为相同后的不同染色数量。
// 每个置换 p 用 p[i] 表示 i 的像，位置从零编号；各位置可独立染色，群大小须在 Z 中可逆。
// 设位置数为 n，时间为 $O(|G|(n+\log n))$。
template <typename Z>
Z burnside(const vector<vector<int>>& group, Z colors) {
    assert(!group.empty());
    Z ret = 0;
    for (const auto& p : group) {
        int n = p.size(), cnt = 0;
        vector<bool> vis(n);
        for (int i = 0; i < n; i++) {
            if (vis[i]) continue;
            cnt++;
            for (int j = i; !vis[j]; j = p[j]) vis[j] = true;
        }
        ret += colors.pow(cnt);
    }
    return ret / Z(group.size());
}

// 传入正位置数 n 和颜色数 colors，返回所有旋转的不动染色数之和，尚未除以 n。
// 试除分解 n 的时间为 $O(\sqrt n)$，其余时间为 $O(\tau(n)(\omega(n)+\log n))$。
template <typename Z>
Z rotation_sum(ll n, Z colors) {
    assert(n >= 1);
    auto fac = prime_factor(n);
    Z ret = 0;
    for (ll d : divisors(fac)) {
        ll phi = d;
        for (auto [p, e] : fac) if (d % p == 0) phi = phi / p * (p - 1);
        ret += Z(phi) * colors.pow(n / d);
    }
    return ret;
}

// 传入正位置数 n 和颜色数 colors，返回只将旋转视为相同的项链数量。
// 要求 n 在 Z 中可逆；各位置可独立染色，颜色有标号。
template <typename Z>
Z necklaces(ll n, Z colors) {
    return rotation_sum(n, colors) / Z(n);
}
// 传入正位置数 n 和颜色数 colors，返回将旋转、翻转均视为相同的手链数量。
// 要求 $2n$ 在 Z 中可逆；各位置可独立染色，颜色有标号。
template <typename Z>
Z bracelets(ll n, Z colors) {
    Z ret = rotation_sum(n, colors);
    if (n & 1) ret += Z(n) * colors.pow(n / 2 + 1);
    else ret += Z(n / 2) * (colors.pow(n / 2 + 1) + colors.pow(n / 2));
    return ret / (Z(n) * 2);
}
\end{lstlisting}

\subsection{Prüfer 编解码（参考）}
\noindent{\small \texttt{Math/Comb/\_prufer.hpp}}\par

\begin{lstlisting}
#pragma once
#include "../head.hpp"

// 传入树的邻接表 g，返回每次删除最小叶子得到的 Prüfer 序列。
// 点从零编号，点数 n 至少为二，返回长度为 $n-2$ 的数组；时间为 $O(n\log n)$。
inline vector<int> prufer_encode(const vector<vector<int>>& g) {
    int n = g.size();
    assert(n >= 2);
    vector<int> deg(n), ret;
    priority_queue<int, vector<int>, greater<int>> q;
    for (int i = 0; i < n; i++) {
        deg[i] = g[i].size();
        if (deg[i] == 1) q.push(i);
    }
    for (int i = 0; i < n - 2; i++) {
        int u = q.top();
        q.pop();
        deg[u] = 0;
        for (int v : g[u]) {
            if (deg[v] == 0) continue;
            ret.push_back(v);
            if (--deg[v] == 1) q.push(v);
            break;
        }
    }
    return ret;
}

// 传入 Prüfer 序列 a，返回对应树的无向边列表，共 $n-1$ 条边，其中 $n=|a|+2$。
// 序列元素须在 $[0,n)$ 内；空序列对应两个点的树，时间为 $O(n\log n)$。
inline vector<pair<int, int>> prufer_decode(const vector<int>& a) {
    int n = a.size() + 2;
    vector<int> deg(n, 1);
    for (int x : a) {
        assert(0 <= x && x < n);
        deg[x]++;
    }
    priority_queue<int, vector<int>, greater<int>> q;
    for (int i = 0; i < n; i++) if (deg[i] == 1) q.push(i);
    vector<pair<int, int>> ret;
    for (int x : a) {
        int u = q.top();
        q.pop();
        ret.emplace_back(u, x);
        if (--deg[x] == 1) q.push(x);
    }
    int u = q.top();
    q.pop();
    ret.emplace_back(u, q.top());
    return ret;
}
\end{lstlisting}

\subsection{FFT}
\noindent{\small \texttt{Math/Poly/FFT.hpp}}\par

\begin{lstlisting}
#pragma once
#include "../head.hpp"

// ok
namespace FFT {

using Real = double;
using cd = complex<Real>;

inline constexpr Real PI = 3.14159265358979323846;
inline int base = 1;
inline vector<int> rev = {0, 1};
inline vector<cd> roots = {cd(0, 0), cd(1, 0)};

// O(2^nbase), 单位根及位逆序缓存
inline void ensure_base(int nbase) {
    if (nbase <= base) return;

    rev.resize(1 << nbase);
    roots.resize(1 << nbase);

    for (int i = 0; i < (1 << nbase); i++) {
        rev[i] = (rev[i >> 1] >> 1)
               | ((i & 1) << (nbase - 1));
    }

    while (base < nbase) {
        Real angle = 2 * PI / (1 << (base + 1));
        for (int i = 1 << (base - 1); i < (1 << base); i++) {
            roots[i << 1] = roots[i];
            Real ang = angle * (2 * i + 1 - (1 << base));
            roots[i << 1 | 1] = cd(cos(ang), sin(ang));
        }
        base++;
    }
}

// O(n log n), n 必须为 2 的幂
inline void fft(vector<cd>& a, int n) {
    int zeros = __builtin_ctz((unsigned)n);
    ensure_base(zeros);
    int shift = base - zeros;

    for (int i = 0; i < n; i++) {
        int j = rev[i] >> shift;
        if (i < j) swap(a[i], a[j]);
    }

    for (int k = 1; k < n; k <<= 1) {
        for (int i = 0; i < n; i += k << 1) {
            for (int j = 0; j < k; j++) {
                cd z = a[i + j + k] * roots[j + k];
                a[i + j + k] = a[i + j] - z;
                a[i + j] += z;
            }
        }
    }
}

inline void ifft(vector<cd>& a, int n) {
    assert(n > 0 && (n & (n - 1)) == 0 && (int)a.size() >= n);

    reverse(a.begin() + 1, a.begin() + n);
    fft(a, n);
    for (int i = 0; i < n; i++) a[i] /= n;
}

// O(L log L), L = ceil_pow2(n + m - 1)
// 支持 vector<double> 和 vector<ll>
template <class T>
vector<T> convolution( const vector<T>& a, const vector<T>& b, int need = -1) {
    static_assert(is_same_v<T, Real> || is_same_v<T, ll>);

    if (a.empty() || b.empty()) return {};

    int full = (int)a.size() + (int)b.size() - 1;
    if (need < 0 || need > full) need = full;
    if (need == 0) return {};

    int na = min<int>(a.size(), need);
    int nb = min<int>(b.size(), need);

    vector<T> c(need);
    if (1LL * na * nb <= 30000) {
        for (int i = 0; i < na; i++) {
            for (int j = 0; j < nb && i + j < need; j++) {
                c[i + j] += a[i] * b[j];
            }
        }
        return c;
    }

    int n = 1;
    while (n < na + nb - 1) n <<= 1;

    vector<cd> fa(n), fb(n);
    for (int i = 0; i < na; i++) fa[i] = (Real)a[i];
    for (int i = 0; i < nb; i++) fb[i] = (Real)b[i];

    fft(fa, n);
    fft(fb, n);
    for (int i = 0; i < n; i++) fa[i] *= fb[i];
    ifft(fa, n);

    for (int i = 0; i < need; i++) {
        if constexpr (is_same_v<T, ll>) {
            c[i] = llround(fa[i].real());
        } else {
            c[i] = fa[i].real();
        }
    }
    return c;
}

} // namespace FFT
\end{lstlisting}

\subsection{NTT 与原根}
\noindent{\small \texttt{Math/Poly/NTT.hpp}}\par

\begin{lstlisting}
#pragma once
#include "../head.hpp"
#include "../NT/factor.hpp"
#include "../NT/tools.hpp"
#include "../NT/modint.hpp"

// ok

// 依赖 modint, prime_factor, mul, qkp

// 返回 m 的最小原根，不存在则返回 -1
// 约O(n^0.25)
inline ll find_primitive_root(ll m) {
    if (m == 2) return 1;
    if (m <= 1) return -1;

    auto fac = prime_factor(m);

    // 原根存在当且仅当 m = 2, 4, p^k, 2p^k (p 为奇素数)
    if (fac[0].first == 2) {
        if (fac.size() == 1) {
            if (fac[0].second != 2) return -1;
        } else {
            if (fac.size() != 2 || fac[0].second != 1) return -1;
        }
    } else if (fac.size() != 1) {
        return -1;
    }

    ll phi = m;
    for (auto [p, cnt] : fac) {
        phi = phi / p * (p - 1);
    }

    auto pf = prime_factor(phi);
    for (ll g = 2; g < m; g++) {
        if (gcd(g, m) != 1) continue;
        bool ok = true;
        for (auto [p, cnt] : pf) {
            if (qkp(g, phi / p, m) == 1) {
                ok = false;
                break;
            }
        }
        if (ok) return g;
    }
    return -1;
}

// O(n log n)
template <class Z, int G = 3>
void ntt(vector<Z>& a, bool inv = false) {
    int n = a.size();

    static vector<int> rev;
    if ((int)rev.size() != n) {
        rev.assign(n, 0);
        for (int i = 1; i < n; i++) {
            rev[i] = (rev[i >> 1] >> 1)
                   | ((i & 1) ? n >> 1 : 0);
        }
    }

    // roots[len + j] = w_(2*len)^j
    static vector<Z> roots{0, 1};
    if ((int)roots.size() < n) {
        int k = __builtin_ctz((unsigned)roots.size());
        roots.resize(n);
        while ((1 << k) < n) {
            Z w = Z(G).pow(
                (Z::mod() - 1) / (1LL << (k + 1))
            );
            for (int i = 1 << (k - 1); i < (1 << k); i++) {
                roots[i << 1] = roots[i];
                roots[i << 1 | 1] = roots[i] * w;
            }
            k++;
        }
    }

    for (int i = 0; i < n; i++) {
        if (i < rev[i]) swap(a[i], a[rev[i]]);
    }

    for (int len = 1; len < n; len <<= 1) {
        for (int i = 0; i < n; i += len << 1) {
            for (int j = 0; j < len; j++) {
                Z x = a[i + j];
                Z y = a[i + j + len] * roots[len + j];
                a[i + j] = x + y;
                a[i + j + len] = x - y;
            }
        }
    }

    if (inv) {
        reverse(a.begin() + 1, a.end());
        Z in = Z(n).inv();
        for (auto& x : a) x *= in;
    }
}

// O(L log L), L = ceil_pow2(n + m - 1)
template <class Z, int G = 3>
vector<Z> convolution(vector<Z> a, vector<Z> b, int need = -1) {
    if (a.empty() || b.empty()) return {};

    int full = (int)a.size() + (int)b.size() - 1;
    if (need == -1 || need > full) need = full;
    if (need <= 0) return {};

    int n = min<int>(a.size(), need);
    int m = min<int>(b.size(), need);
    a.resize(n);
    b.resize(m);

    if (min(n, m) <= 60) {
        vector<Z> c(need);
        for (int i = 0; i < n; i++) {
            for (int j = 0; j < m && i + j < need; j++) {
                c[i + j] += a[i] * b[j];
            }
        }
        return c;
    }

    int len = 1;
    while (len < n + m - 1) len <<= 1;
    assert((Z::mod() - 1) % len == 0);

    a.resize(len);
    b.resize(len);
    ntt<Z, G>(a);
    ntt<Z, G>(b);

    for (int i = 0; i < len; i++) a[i] *= b[i];

    ntt<Z, G>(a, true);
    a.resize(need);
    return a;
}
\end{lstlisting}

\subsection{三模 NTT 与精确整数卷积}
\noindent{\small \texttt{Math/Poly/MTT.hpp}}\par

\begin{lstlisting}
#pragma once
#include "../head.hpp"
#include "../NT/tools.hpp"
#include "../NT/modint.hpp"
#include "NTT.hpp"
// ok

// 三模 NTT + CRT
// O(L log L), L = ceil_pow2(n + m - 1)
// Require: mod > 0, NTT length <= 2^24
// 对 mod <= 1e9+7，在长度限制内保证正确
// 常用原根为3NTT友好模数:
// 167772161   = 5   * 2^25 + 1
// 469762049   = 7   * 2^26 + 1
// 998244353   = 119 * 2^23 + 1
// 1224736769  = 73  * 2^24 + 1
// 2130706433  = 127 * 2^24 + 1
// 1004535809  = 479 * 2^21 + 1
inline vector<ll> mtt(const vector<ll>& a, const vector<ll>& b, ll mod, int need = -1) {
    if (a.empty() || b.empty()) return {};

    int full = (int)a.size() + b.size() - 1;
    if (need == -1 || need > full) need = full;
    if (need <= 0) return {};

    int n = min<int>(a.size(), need);
    int m = min<int>(b.size(), need);

    if (1LL * n * m <= 30000) {
        vector<ll> c(need);
        for (int i = 0; i < n; i++) {
            ll x = norm(a[i], mod);
            for (int j = 0; j < m && i + j < need; j++) {
                ll y = norm(b[j], mod);
                c[i + j] = (c[i + j] + (i128)x * y) % mod;
            }
        }
        return c;
    }

    constexpr ll M1 = 167772161;
    constexpr ll M2 = 469762049;
    constexpr ll M3 = 1224736769;

    int len = 1;
    while (len < n + m - 1) len <<= 1;
    assert(len <= (1 << 24));

    using Z1 = ModInt<M1>;
    using Z2 = ModInt<M2>;
    using Z3 = ModInt<M3>;

    vector<Z1> a1(n), b1(m);
    vector<Z2> a2(n), b2(m);
    vector<Z3> a3(n), b3(m);

    for (int i = 0; i < n; i++) {
        ll x = norm(a[i], mod);
        a1[i] = x;
        a2[i] = x;
        a3[i] = x;
    }
    for (int i = 0; i < m; i++) {
        ll x = norm(b[i], mod);
        b1[i] = x;
        b2[i] = x;
        b3[i] = x;
    }

    auto c1 = convolution(move(a1), move(b1), need);
    auto c2 = convolution(move(a2), move(b2), need);
    auto c3 = convolution(move(a3), move(b3), need);

    static const ll inv12 = qkp(M1, M2 - 2, M2);
    static const ll inv123 = qkp(M1 * M2 % M3, M3 - 2, M3);

    ll m12 = M1 * M2 % mod;
    vector<ll> ans(need);

    for (int i = 0; i < need; i++) {
        ll x1 = c1[i].val;
        ll x2 = c2[i].val;
        ll x3 = c3[i].val;

        ll t1 = (x2 - x1 % M2 + M2) % M2 * inv12 % M2;
        ll x12 = x1 + t1 * M1;
        ll t2 = (x3 - x12 % M3 + M3) % M3 * inv123 % M3;

        ans[i] = (x12 % mod + (i128)t2 * m12) % mod;
    }
    return ans;
}

// 精确整数卷积, 结果在ll范围内
// Require: NTT length <= 2^24
template <int P>
vector<ModInt<P>> conv_ll_mod(const vector<ll>& a, const vector<ll>& b, int need) {
    using Z = ModInt<P>;
    vector<Z> x(a.begin(), a.end());
    vector<Z> y(b.begin(), b.end());
    return convolution(move(x), move(y), need);
}

inline vector<ll> convolution_ll(const vector<ll>& a, const vector<ll>& b, int need = -1) {
    if (a.empty() || b.empty()) return {};

    int full = (int)a.size() + (int)b.size() - 1;
    if (need == -1 || need > full) need = full;
    if (need <= 0) return {};

    constexpr ll P1 = 167772161;
    constexpr ll P2 = 469762049;
    constexpr ll P3 = 1224736769;
    constexpr i128 P = (i128)P1 * P2 * P3;

    auto c1 = conv_ll_mod<P1>(a, b, need);
    auto c2 = conv_ll_mod<P2>(a, b, need);
    auto c3 = conv_ll_mod<P3>(a, b, need);

    static const ll inv12 = qkp(P1, P2 - 2, P2);
    static const ll inv123 = qkp(P1 * P2 % P3, P3 - 2, P3);

    vector<ll> ans(need);
    for (int i = 0; i < need; i++) {
        ll x1 = c1[i].val;
        ll x2 = c2[i].val;
        ll x3 = c3[i].val;

        ll t1 = (x2 - x1 + P2) % P2 * inv12 % P2;
        ll x12 = x1 + P1 * t1;
        ll t2 = (x3 - x12 % P3 + P3) % P3 * inv123 % P3;

        i128 x = x12 + (i128)P1 * P2 * t2;
        if (x > P / 2) x -= P;

        ans[i] = (ll)x;
    }
    return ans;
}
\end{lstlisting}

\subsection{多模 NTT 与 Garner 合并}
\noindent{\small \texttt{Math/Poly/MTT\_large.hpp}}\par

\begin{lstlisting}
#pragma once
#include "../head.hpp"
#include "../NT/tools.hpp"
#include "../NT/modint.hpp"
#include "NTT.hpp"
// ok maybe

// 单个 NTT 模数配置
template <int P, int G = 3>
struct NTTPrime {
    static constexpr int mod = P, root = G;
};

// 单模卷积，返回原始余数
template <int P, int G>
vector<u32> mtt_residue(
    const vector<ll>& a, const vector<ll>& b, int need
) {
    using Z = ModInt<P>;
    vector<Z> x(a.size()), y(b.size());

    for (int i = 0; i < (int)a.size(); i++) x[i] = a[i];
    for (int i = 0; i < (int)b.size(); i++) y[i] = b[i];

    auto c = convolution<Z, G>(move(x), move(y), need);
    vector<u32> ret(need);
    for (int i = 0; i < need; i++) ret[i] = c[i].val;
    return ret;
}

// K 模 NTT + Garner
// Require: mod > 0, primes are distinct NTT primes
// Require: integer convolution coefficients < product(primes)
template <class... Cfg>
vector<ll> mtt_k(
    const vector<ll>& a, const vector<ll>& b,
    ll mod, int need = -1
) {
    constexpr int K = sizeof...(Cfg);
    static_assert(K > 0);
    constexpr array<ll, K> ps = {Cfg::mod...};

    assert(mod > 0);
    if (a.empty() || b.empty()) return {};

    int full = (int)a.size() + (int)b.size() - 1;
    if (need == -1 || need > full) need = full;
    if (need <= 0) return {};
    if (mod == 1) return vector<ll>(need, 0);

    int n = min<int>(a.size(), need);
    int m = min<int>(b.size(), need);

    vector<ll> aa(n), bb(m);
    for (int i = 0; i < n; i++) aa[i] = norm(a[i], mod);
    for (int i = 0; i < m; i++) bb[i] = norm(b[i], mod);

    if (1LL * n * m <= 30000) {
        vector<ll> c(need);
        for (int i = 0; i < n; i++) {
            for (int j = 0; j < m && i + j < need; j++) {
                c[i + j] = (c[i + j] + (i128)aa[i] * bb[j]) % mod;
            }
        }
        return c;
    }

    int len = 1;
    while (len < n + m - 1) len <<= 1;
    assert(( ((Cfg::mod - 1) % len == 0) && ... ));

    // 分别在 K 个模数下进行卷积
    array<vector<u32>, K> rs = {
        mtt_residue<Cfg::mod, Cfg::root>(aa, bb, need)...
    };

    // invs[j][i] = inverse(ps[j]) mod ps[i]
    static const auto invs = [] {
        array<array<ll, K>, K> v{};
        constexpr array<ll, K> p = {Cfg::mod...};
        for (int i = 0; i < K; i++) {
            for (int j = 0; j < i; j++) {
                v[j][i] = qkp(p[j] % p[i], p[i] - 2, p[i]);
            }
        }
        return v;
    }();

    // pref[i] = product(ps[0...i-1]) % mod
    array<ll, K> pref{};
    pref[0] = 1 % mod;
    for (int i = 1; i < K; i++) {
        pref[i] = (i128)pref[i - 1] * ps[i - 1] % mod;
    }

    vector<ll> ans(need);

    for (int i = 0; i < need; i++) {
        array<ll, K> digit{};

        // Garner: 计算混合进制系数
        for (int j = 0; j < K; j++) {
            ll x = rs[j][i];
            for (int t = 0; t < j; t++) {
                x = (x - digit[t]) % ps[j];
                if (x < 0) x += ps[j];
                x = (i128)x * invs[t][j] % ps[j];
            }
            digit[j] = x;
        }

        // 直接在目标模数下还原
        ll x = 0;
        for (int j = 0; j < K; j++) {
            x = (x + (i128)digit[j] * pref[j]) % mod;
        }
        ans[i] = x;
    }
    return ans;
}

// 五模配置
using P0 = NTTPrime<167772161>;
using P1 = NTTPrime<469762049>;
using P2 = NTTPrime<1224736769>;
using P3 = NTTPrime<2130706433>;
using P4 = NTTPrime<2013265921, 31>;

inline vector<ll> mtt(
    const vector<ll>& a, const vector<ll>& b,
    ll mod, int need = -1
) {
    return mtt_k<P0, P1, P2, P3, P4>(a, b, mod, need);
}
\end{lstlisting}

\subsection{按位卷积}
\noindent{\small \texttt{Math/Poly/FWHT.hpp}}\par

\begin{lstlisting}
#pragma once
#include "../head.hpp"

// ok
template <typename Z>
void fwt_or(vector<Z>& a, bool invert = 0) {
    int n = a.size();
    for (int len = 1; len < n; len <<= 1) {
        for (int i = 0; i < n; i += len << 1) {
            for (int j = 0; j < len; j++) {
                if (!invert) {
                    a[i + j + len] += a[i + j];
                } else {
                    a[i + j + len] -= a[i + j];
                }
            }
        }
    }
}
template <typename Z>
void fwt_and(vector<Z>& a, bool invert = 0) {
    int n = a.size();
    for (int len = 1; len < n; len <<= 1) {
        for (int i = 0; i < n; i += len << 1) {
            for (int j = 0; j < len; j++) {
                if (!invert) {
                    a[i + j] += a[i + j + len];
                } else {
                    a[i + j] -= a[i + j + len];
                }
            }
        }
    }
}

template <typename Z>
void fwt_xor(vector<Z>& a, bool invert = 0) {
    int n = a.size();
    for (int len = 1; len < n; len <<= 1) {
        for (int i = 0; i < n; i += len << 1) {
            for (int j = 0; j < len; j++) {
                Z x = a[i + j];
                Z y = a[i + j + len];
                a[i + j] = x + y;
                a[i + j + len] = x - y;
            }
        }
    }

    if (invert) {
        Z invn = Z(n).inv();
        for (int i = 0; i < n; i++) {
            a[i] *= invn;
        }
    }
}
template <typename Z>
vector<Z> convolution_or(vector<Z> a, vector<Z> b) {
    int n = a.size();
    assert((int)b.size() == n);

    fwt_or(a);
    fwt_or(b);

    for (int i = 0; i < n; i++) {
        a[i] *= b[i];
    }

    fwt_or(a, 1);
    return a;
}

template <typename Z>
vector<Z> convolution_and(vector<Z> a, vector<Z> b) {
    int n = a.size();
    assert((int)b.size() == n);

    fwt_and(a);
    fwt_and(b);

    for (int i = 0; i < n; i++) {
        a[i] *= b[i];
    }

    fwt_and(a, 1);
    return a;
}

template <typename Z>
vector<Z> convolution_xor(vector<Z> a, vector<Z> b) {
    int n = a.size();
    assert((int)b.size() == n);

    fwt_xor(a);
    fwt_xor(b);

    for (int i = 0; i < n; i++) {
        a[i] *= b[i];
    }

    fwt_xor(a, 1);
    return a;
}
\end{lstlisting}

\subsection{多项式与形式幂级数运算}
\noindent{\small \texttt{Math/Poly/poly\_major.hpp}}\par

\begin{lstlisting}
#include "../head.hpp"
#include "NTT.hpp"

template<typename Z>
vector<Z> inv(const vector<Z>& a, int n) {
    if (n == 0) return {};
    // assert(a[0] != 0);
    // B = B'(2 - AB')
    int m = 1;
    vector<Z> b(1, a[0].inv());
    while (m < n) {
        int len = min(m * 2, n);
        vector<Z> A(min(len, (int)a.size()));
        for (int i = 0; i < (int)A.size(); i++) {
            A[i] = a[i];
        }
        vector<Z> B = b;
        auto C = convolution(A, B);
        C.resize(len);
        for (int i = 0; i < len; i++) {
            C[i] = -C[i];
        }
        C[0] += 2;
        b = convolution(B, C);
        b.resize(len);
        m *= 2;
    }
    b.resize(n);
    return b;
}

template <typename Z>
vector<Z> deriv(const vector<Z>& a) {
    int n = a.size();
    vector<Z> ret(max(0, n - 1));
    for (int i = 0; i < n - 1; i++) {
        ret[i] = a[i + 1] * Z(i + 1);
    }
    return ret;
}
template <typename Z>
vector<Z> intgr(const vector<Z>& a) {
    int n = a.size();
    assert(n < Z::mod());

    static vector<Z> invs{0, 1};
    int p = Z::mod();

    for (int i = invs.size(); i <= n; i++) {
        invs.push_back(-Z(p / i) * invs[p % i]);
    }

    vector<Z> ret(n + 1);
    for (int i = 1; i <= n; i++) {
        ret[i] = a[i - 1] * invs[i];
    }
    return ret;
}


template <typename Z>
pair<vector<Z>, vector<Z>> divmod(vector<Z> a, vector<Z> b) {
    assert(!b.empty() && b.back() != Z(0));
    int n = a.size();
    int m = b.size();
    if (n < m) {
        return { vector<Z>{ Z(0) }, a };
    }
    int len = n - m + 1;
    vector<Z> ra = a;
    vector<Z> rb = b;
    reverse(ra.begin(), ra.end());
    reverse(rb.begin(), rb.end());
    ra.resize(len);
    rb.resize(len);
    auto q = convolution(ra, inv(rb, len));
    q.resize(len);
    reverse(q.begin(), q.end());
    auto c = convolution(q, b, m - 1);
    a.resize(m - 1);
    for (int i = 0; i < (int)c.size(); i++) {
        a[i] -= c[i];
    }
    return {q, a};
}

template <typename Z>
vector<Z> ln(const vector<Z>& a, int n) {
    // B = ln(A) = intgr(deriv(A) * inv(A))
    // assert(a[0] == 1);
    vector<Z> A(n);
    for (int i = 0; i < n; i++) {
        if (i < a.size()) {
            A[i] = a[i];
        }
    }
    auto B = intgr(convolution(deriv(A), inv(A, n)));
    B.resize(n);
    return B;
}


template <typename Z>
vector<Z> exp(vector<Z> a, int n) {
    // assert(a[0] == 0);
    int m = 1;
    vector<Z> b(1);
    b[0] = 1;
    while (m < n) {
        // B = B' * (1 - ln(B') + A)
        int len = min(m * 2, n);
        vector<Z> A(len);
        for (int i = 0; i < min((int)a.size(), len); i++) {
            A[i] = a[i];
        }
        vector<Z> B(b);
        B.resize(len);
        auto C = ln(B, len);
        for (int i = 0; i < len; i++) {
            C[i] = A[i] - C[i];
        }
        C[0] += 1;
        b = convolution(B, C);
        b.resize(len);
        m *= 2;
    }
    b.resize(n);
    return b;
}


template <typename Z>
vector<Z> mulT(vector<Z> a, vector<Z> b, int need) {
    vector<Z> ret(need);
    if (a.empty() || b.empty() || need == 0) {
        return ret;
    }
    int n = b.size();
    reverse(b.begin(), b.end());
    auto c = convolution(a, b);
    for (int i = 0; i < need; i++) {
        int p = i + n - 1;
        if (p < (int)c.size()) {
            ret[i] = c[p];
        }
    }
    return ret;
}

template <typename Z>
vector<Z> multipoint_eval(vector<Z> f, vector<Z> x) {
    int m = x.size();
    vector<Z> ans(m);
    if (m == 0) {
        return ans;
    }
    if (f.empty()) {
        return ans;
    }
    int n = max((int)f.size(), m);
    x.resize(n);
    vector<vector<Z>> q(4 * n);
    auto build = [&](auto self, int p, int l, int r) -> void {
        if (r - l == 1) {
            q[p] = {Z(1), -x[l]};
            return;
        }
        int mid = (l + r) / 2;
        self(self, p << 1, l, mid);
        self(self, p << 1 | 1, mid, r);
        q[p] = convolution(q[p << 1], q[p << 1 | 1]);
    };
    auto work = [&](auto self, int p, int l, int r, const vector<Z>& num) -> void {
        if (r - l == 1) {
            if (l < m) {
                ans[l] = num.empty() ? Z(0) : num[0];
            }
            return;
        }
        int mid = (l + r) / 2;
        auto vl = mulT(num, q[p << 1 | 1], mid - l);
        auto vr = mulT(num, q[p << 1], r - mid);
        self(self, p << 1, l, mid, vl);
        self(self, p << 1 | 1, mid, r, vr);
    };
    build(build, 1, 0, n);
    auto root = mulT(f, inv(q[1], n), n);
    work(work, 1, 0, n, root);
    return ans;
}

template <typename Z>
vector<Z> interpolate(vector<Z> x, vector<Z> y) {
    int n = x.size();
    assert((int)y.size() == n);
    if (n == 0) {
        return {};
    }
    vector<vector<Z>> q(4 * n);
    auto build = [&](auto self, int p, int l, int r) -> void {
        if (r - l == 1) {
            // x - x_l
            q[p] = {-x[l], Z(1)};
            return;
        }

        int mid = (l + r) / 2;

        self(self, p << 1, l, mid);
        self(self, p << 1 | 1, mid, r);

        q[p] = convolution(q[p << 1], q[p << 1 | 1]);
    };

    build(build, 1, 0, n);
    auto d = multipoint_eval(deriv(q[1]), x);
    vector<Z> w(n);
    for (int i = 0; i < n; i++) {
        w[i] = y[i] / d[i];
    }

    auto work = [&](auto self, int p, int l, int r) -> vector<Z> {
        if (r - l == 1) {
            return vector<Z>{w[l]};
        }

        int mid = (l + r) / 2;

        auto L = self(self, p << 1, l, mid);
        auto R = self(self, p << 1 | 1, mid, r);

        auto A = convolution(L, q[p << 1 | 1]);
        auto B = convolution(R, q[p << 1]);

        int len = max(A.size(), B.size());
        A.resize(len);
        B.resize(len);

        for (int i = 0; i < len; i++) {
            A[i] += B[i];
        }

        return A;
    };
    auto ans = work(work, 1, 0, n);
    ans.resize(n);
    return ans;
}


template <typename Z>
vector<Z> pow(vector<Z> a, ll b, int n) {
    if (b < 0) {
        return (inv(a, n), -b, n);
    }
    if (n == 0) return {};
    if (b == 0) {
        vector<Z> ret(n);
        ret[0] = 1;
        return ret;
    }
    a.resize(n);
    int t = n;
    for (int i = 0; i < n; i++) {
        if (a[i] != Z(0)) {
            t = i;
            break;
        }
    }
    vector<Z> ans(n);
    if (t == n) {
        return ans;
    }
    if ((__int128)t * b >= n) {
        return ans;
    }
    int s = t * b;
    int m = n - s;
    vector<Z> c(m);
    for (int i = 0; i < m && i + t < n; i++) {
        c[i] = a[i + t] / a[t];
    }
    c = ln(c, m);
    for (int i = 0; i < m; i++) {
        c[i] *= Z(b);
    }
    c = exp(c, m);
    Z l = a[t].pow(b);
    for (int i = 0; i < m; i++) {
        ans[i + s] = c[i] * l;
    }
    return ans;
}

template <typename Z>
vector<Z> sqrt(const vector<Z>& a, int n) {
    // assert(a[0] == 1);
    int m = 1;
    vector<Z> b(1);
    b[0] = 1;

    Z inv2 = Z(2).inv();
    while (m < n) {
        // B = (B' + A / B') / 2
        int len = min(m * 2, n);
        vector<Z> A(len);
        for (int i = 0; i < min((int)a.size(), len); i++) {
            A[i] = a[i];
        }
        b.resize(len);
        vector<Z> B(b);
        b = convolution(A, inv(B, len));
        b.resize(len);
        for (int i = 0; i < len; i++) {
            b[i] += B[i];
        }
        for (int i = 0; i < len; i++) {
            b[i] *= inv2;
        }
        m *= 2;
    }
    b.resize(n);
    return b;
}
\end{lstlisting}

\subsection{BM 与 Kitamasa 线性递推（参考）}
\noindent{\small \texttt{Math/Poly/\_linear\_recurrence.hpp}}\par

\begin{lstlisting}
#pragma once
#include "../head.hpp"
#include "../NT/modint.hpp"

// 传入数列前缀 a，返回与它一致的最短递推系数 c，c[0] 对应最近一项；全零前缀返回空数组。
// Z 必须是域元素；n 为前缀长度，时间为 $O(n^2)$。该结果是否对后续成立，仍需题目结构保证。
template <typename Z>
vector<Z> berlekamp_massey(const vector<Z>& a) {
    vector<Z> c{1}, b{1};
    int len = 0, step = 1;
    Z last = 1;
    for (int i = 0; i < (int)a.size(); i++) {
        Z d = a[i];
        for (int j = 1; j <= len; j++) d += c[j] * a[i - j];
        if (d == Z(0)) {
            step++;
            continue;
        }
        auto old = c;
        Z t = d / last;
        c.resize(max(c.size(), b.size() + step));
        for (int j = 0; j < (int)b.size(); j++) c[j + step] -= t * b[j];
        if (2 * len <= i) {
            len = i + 1 - len;
            b = move(old);
            last = d;
            step = 1;
        } else {
            step++;
        }
    }
    c.resize(len + 1);
    c.erase(c.begin());
    for (auto& x : c) x = -x;
    return c;
}

// 传入已知前缀 a、递推系数 c 和非负下标 n，返回第 n 项；c[0] 对应最近一项。
// a 至少含前 k 项，k 为 c 的长度；空递推的后续项为零。Kitamasa 时间为 $O(k^2\log n)$，支持合数模。
template <typename Z>
Z linear_nth(const vector<Z>& a, const vector<Z>& c, ll n) {
    assert(n >= 0 && a.size() >= c.size());
    if (n < (ll)a.size()) return a[n];
    int k = c.size();
    if (k == 0) return 0;
    // 传入两个次数小于 k 的多项式系数数组 a、b，返回乘积对递推特征多项式取模后的系数。
    auto combine = [&](const vector<Z>& a, const vector<Z>& b) {
        vector<Z> ret(2 * k - 1);
        for (int i = 0; i < k; i++) {
            for (int j = 0; j < k; j++) ret[i + j] += a[i] * b[j];
        }
        for (int i = 2 * k - 2; i >= k; i--) {
            for (int j = 0; j < k; j++) ret[i - 1 - j] += ret[i] * c[j];
        }
        ret.resize(k);
        return ret;
    };
    vector<Z> ret(k), mul(k);
    ret[0] = 1;
    if (k == 1) mul[0] = c[0];
    else mul[1] = 1;
    for (; n; n >>= 1, mul = combine(mul, mul)) {
        if (n & 1) ret = combine(ret, mul);
    }
    Z ans = 0;
    for (int i = 0; i < k; i++) ans += ret[i] * a[i];
    return ans;
}

// 传入数列前缀 a 和非负下标 n，先用 BM 恢复递推，再返回据此推算的第 n 项。
// Z 必须是域元素；递推在前缀之外是否成立，须由调用者根据题目保证。
template <typename Z>
Z linear_guess(const vector<Z>& a, ll n) {
    return linear_nth(a, berlekamp_massey(a), n);
}
\end{lstlisting}

\subsection{Bostan-Mori 与快速线性递推（参考）}
\noindent{\small \texttt{Math/Poly/\_bostan\_mori.hpp}}\par

\begin{lstlisting}
#pragma once
#include "NTT.hpp"

// 传入按升幂存储的分子 p、分母 q 和非负下标 n，返回形式幂级数 $P(x)/Q(x)$ 的第 n 项系数。
// 要求 q 非空、q[0] 可逆，并满足现有 NTT 的条件；当 $\deg P<\deg Q=k$ 时，时间为 $O(k\log k\log n)$。
template <typename Z>
Z bostan_mori(vector<Z> p, vector<Z> q, ll n) {
    assert(n >= 0 && !q.empty() && q[0] != Z(0));
    while (n && !p.empty()) {
        auto b = q;
        for (int i = 1; i < (int)b.size(); i += 2) b[i] = -b[i];
        auto a = convolution(p, b), c = convolution(q, b);
        p.clear();
        q.clear();
        for (int i = n & 1; i < (int)a.size(); i += 2) p.push_back(a[i]);
        for (int i = 0; i < (int)c.size(); i += 2) q.push_back(c[i]);
        n >>= 1;
    }
    return p.empty() ? Z(0) : p[0] / q[0];
}

// 传入已知前缀 a、递推系数 c 和非负下标 n，返回第 n 项；c[0] 对应最近一项。
// a 至少包含前 k 项，其中 k 是 c 的长度；空递推的后续项为零。使用现有 NTT 加速。
template <typename Z>
Z linear_nth_ntt(const vector<Z>& a, const vector<Z>& c, ll n) {
    assert(n >= 0 && a.size() >= c.size());
    if (n < (ll)a.size()) return a[n];
    int k = c.size();
    if (k == 0) return 0;
    vector<Z> q(k + 1), b(a.begin(), a.begin() + k);
    q[0] = 1;
    for (int i = 0; i < k; i++) q[i + 1] = -c[i];
    auto p = convolution(b, q, k);
    return bostan_mori(move(p), move(q), n);
}
\end{lstlisting}

\subsection{连续点插值与幂和（参考）}
\noindent{\small \texttt{Math/Poly/\_lagrange.hpp}}\par

\begin{lstlisting}
#pragma once
#include "../head.hpp"
#include "../NT/modint.hpp"

// 传入连续点函数值 y 和求值位置 x，返回多项式在 x 处的值；y[i] 是在横坐标 i 处的值。
// 若 y 的长度为 $k+1$，要求 $\deg f\le k<P$，Z 为 ModInt<P> 且 P 为质数；空 y 返回零，时间为 $O(k)$。
template <typename Z>
Z lagrange(const vector<Z>& y, ll x) {
    int n = y.size();
    if (n == 0) return 0;
    assert(n <= Z::mod());
    Z t = x;
    if (t.val < (u32)n) return y[t.val];
    vector<Z> pre(n + 1, 1), suf(n + 1, 1), factinv(n);
    Z fact = 1;
    for (int i = 1; i < n; i++) fact *= i;
    factinv[n - 1] = fact.inv();
    for (int i = n - 1; i > 0; i--) factinv[i - 1] = factinv[i] * i;
    for (int i = 0; i < n; i++) pre[i + 1] = pre[i] * (t - Z(i));
    for (int i = n - 1; i >= 0; i--) suf[i] = suf[i + 1] * (t - Z(i));
    Z ret = 0;
    for (int i = 0; i < n; i++) {
        Z val = y[i] * pre[i] * suf[i + 1] * factinv[i] * factinv[n - 1 - i];
        ret += (n - 1 - i) & 1 ? -val : val;
    }
    return ret;
}

// 传入非负整数 n、k，返回从一到 n 的 k 次幂之和，结果按 Z 的模数计算。
// Z 为质数模 ModInt<P>，要求 $k+1<P$；时间为 $O(k\log(k+1))$。
template <typename Z = mint>
Z power_sum(ll n, int k) {
    assert(n >= 0 && k >= 0 && (ll)k + 1 < Z::mod());
    vector<Z> y(k + 2);
    for (int i = 1; i <= k + 1; i++) y[i] = y[i - 1] + Z(i).pow(k);
    return lagrange(y, n);
}
\end{lstlisting}

\subsection{Taylor shift 多项式平移（参考）}
\noindent{\small \texttt{Math/Poly/\_taylor\_shift.hpp}}\par

\begin{lstlisting}
#pragma once
#include "NTT.hpp"

// 传入多项式的升幂系数 a 和常数 c，返回 $f(x+c)$ 的升幂系数，长度不变；空输入返回空数组。
// Z 为质数模类型，要求数组长度不超过模数，并满足现有 NTT 的条件；时间为 $O(n\log n)$。
template <typename Z>
vector<Z> taylor_shift(const vector<Z>& a, Z c) {
    int n = a.size();
    if (n == 0) return {};
    assert(n <= Z::mod());
    vector<Z> fact(n, 1), factinv(n, 1), b(n), d(n);
    for (int i = 1; i < n; i++) fact[i] = fact[i - 1] * i;
    factinv[n - 1] = fact[n - 1].inv();
    for (int i = n - 1; i > 0; i--) factinv[i - 1] = factinv[i] * i;
    Z mul = 1;
    for (int i = 0; i < n; i++) {
        b[n - 1 - i] = a[i] * fact[i];
        d[i] = mul * factinv[i];
        mul *= c;
    }
    auto ret = convolution(b, d, n);
    reverse(ret.begin(), ret.end());
    for (int i = 0; i < n; i++) ret[i] *= factinv[i];
    return ret;
}
\end{lstlisting}

\subsection{子集卷积（参考）}
\noindent{\small \texttt{Math/Poly/\_subset\_convolution.hpp}}\par

\begin{lstlisting}
#pragma once
#include "FWHT.hpp"

// 传入按集合掩码存储、长度同为 $2^k$ 的非空数组 a、b，返回只合并不相交集合的卷积。
// 结果为 $c_S=\sum_{T\subseteq S}a_T b_{S\setminus T}$；时间为 $O(k^2 2^k)$。
// 主要数组存放约 $2(k+1)2^k$ 个 Z；当 $k=20$ 且 Z 为 mint 时，约占 $168\,\mathrm{MiB}$。
template <typename Z>
vector<Z> subset_convolution(const vector<Z>& a, const vector<Z>& b) {
    int n = a.size();
    assert(n > 0 && (n & (n - 1)) == 0 && a.size() == b.size());
    int k = __builtin_ctz((unsigned)n);
    vector<int> cnt(n);
    vector<vector<Z>> A(k + 1, vector<Z>(n)), B = A;
    for (int s = 0; s < n; s++) {
        cnt[s] = __builtin_popcount((unsigned)s);
        A[cnt[s]][s] = a[s];
        B[cnt[s]][s] = b[s];
    }
    for (int i = 0; i <= k; i++) {
        fwt_or(A[i]);
        fwt_or(B[i]);
    }
    vector<Z> ret(n);
    for (int i = 0; i <= k; i++) {
        vector<Z> c(n);
        for (int j = 0; j <= i; j++) {
            for (int s = 0; s < n; s++) c[s] += A[j][s] * B[i - j][s];
        }
        fwt_or(c, true);
        for (int s = 0; s < n; s++) if (cnt[s] == i) ret[s] = c[s];
    }
    return ret;
}
\end{lstlisting}

\subsection{高斯消元}
\noindent{\small \texttt{Math/Alg/gauss.hpp}}\par

\begin{lstlisting}
#pragma once
#include "../head.hpp"
#include "../NT/modint.hpp"

// O(n * m * min(n, m))
// 求解 Ax = b，a 为 n 行 m+1 列增广矩阵
// 返回 {rank, x}，无解时 rank = -1
// rank = m 时唯一解，否则有自由变量（取 0）
template <typename Z>
pair<int, vector<Z>> gauss(vector<vector<Z>> a, double eps = 1e-9) {
    int n = a.size(), m = a[0].size() - 1;
    vector<int> pos;
    int rank = 0;

    auto zero = [&](Z x) {
        if constexpr (is_floating_point_v<Z>) {
            return abs(x) <= eps;
        } else {
            return x == Z(0);
        }
    };

    for (int col = 0; col < m && rank < n; col++) {
        int pivot = rank;
        if constexpr (is_floating_point_v<Z>) {
            for (int i = rank + 1; i < n; i++) {
                if (abs(a[i][col]) > abs(a[pivot][col])) {
                    pivot = i;
                }
            }
            if (zero(a[pivot][col])) continue;
        } else {
            while (pivot < n && zero(a[pivot][col])) {
                pivot++;
            }
            if (pivot == n) continue;
        }

        swap(a[pivot], a[rank]);
        Z iv = Z(1) / a[rank][col];
        for (int j = col; j <= m; j++) {
            a[rank][j] *= iv;
        }

        for (int i = rank + 1; i < n; i++) {
            Z t = a[i][col];
            if (t == Z(0)) continue;
            for (int j = col; j <= m; j++) {
                a[i][j] -= t * a[rank][j];
            }
        }
        pos.push_back(col);
        rank++;
    }

    for (int i = rank; i < n; i++) {
        if (!zero(a[i][m])) {
            return {-1, {}};
        }
    }

    vector<Z> x(m);
    for (int i = rank - 1; i >= 0; i--) {
        int col = pos[i];
        Z val = a[i][m];
        for (int j = col + 1; j < m; j++) {
            val -= a[i][j] * x[j];
        }
        x[col] = val;
    }
    
    // 基础解系
    // vector<bool> is_pivot(m);
    // for (int col : pos) {
    //     is_pivot[col] = true;
    // }

    // vector<vector<Z>> basis;
    // for (int j = 0; j < m; j++) {
    //     if (is_pivot[j]) continue;

    //     vector<Z> v(m);
    //     v[j] = 1;

    //     for (int i = rank - 1; i >= 0; i--) {
    //         int col = pos[i];
    //         for (int k = col + 1; k < m; k++) {
    //             v[col] -= a[i][k] * v[k];
    //         }
    //     }
    //     basis.push_back(v);
    // }
    // return { rank, x, basis };
    return { rank, x };
}
\end{lstlisting}

\subsection{矩阵与矩阵快速幂（参考）}
\noindent{\small \texttt{Math/Alg/\_Matrix.hpp}}\par

\begin{lstlisting}
#pragma once
#include "../head.hpp"
#include "../NT/modint.hpp"

// 存储矩阵并提供加减、乘法和非负整数次幂；Z 是元素类型。
// 可使用 mint、ll 或浮点类型；整数运算须保证结果不溢出。
template <typename Z>
struct Matrix {
    int n, m;
    vector<vector<Z>> a;
    // 传入非负行数 n、列数 m，构造一个全零矩阵。
    Matrix(int n, int m) : n(n), m(m), a(n, vector<Z>(m)) {}
    // 传入各行等长的二维数组 a，以其内容构造矩阵；空数组表示空矩阵。
    Matrix(vector<vector<Z>> a) : n(a.size()), m(n ? a[0].size() : 0), a(move(a)) {}
    // 传入行下标 i，返回该行的引用，可直接读写元素。
    vector<Z>& operator[](int i) { return a[i]; }
    // 传入行下标 i，返回该行的只读引用。
    const vector<Z>& operator[](int i) const { return a[i]; }

    // 传入阶数 n，返回 n 阶单位矩阵。
    static Matrix identity(int n) {
        Matrix ret(n, n);
        for (int i = 0; i < n; i++) ret[i][i] = 1;
        return ret;
    }
    // 传入同型矩阵 b，返回逐项相加的结果，不修改原矩阵。
    Matrix operator+(const Matrix& b) const {
        assert(n == b.n && m == b.m);
        Matrix ret = *this;
        for (int i = 0; i < n; i++) {
            for (int j = 0; j < m; j++) ret[i][j] += b[i][j];
        }
        return ret;
    }
    // 传入同型矩阵 b，返回逐项相减的结果，不修改原矩阵。
    Matrix operator-(const Matrix& b) const {
        assert(n == b.n && m == b.m);
        Matrix ret = *this;
        for (int i = 0; i < n; i++) {
            for (int j = 0; j < m; j++) ret[i][j] -= b[i][j];
        }
        return ret;
    }
    
    // 传入行数等于当前列数的矩阵 b，返回矩阵乘积。
    // 按行访问并跳过零系数；时间为 $O(nmq)$，q 是 b 的列数。
    Matrix operator*(const Matrix& b) const {
        assert(m == b.n);
        Matrix ret(n, b.m);
        for (int i = 0; i < n; i++) {
            for (int k = 0; k < m; k++) {
                Z t = a[i][k];
                if (t == Z(0)) continue;
                for (int j = 0; j < b.m; j++) ret[i][j] += t * b[k][j];
            }
        }
        return ret;
    }
    // 传入长度为 m 的列向量 b，返回矩阵与列向量的乘积。
    vector<Z> operator*(const vector<Z>& b) const {
        assert(m == (int)b.size());
        vector<Z> ret(n);
        for (int i = 0; i < n; i++) {
            for (int j = 0; j < m; j++) ret[i] += a[i][j] * b[j];
        }
        return ret;
    }
    
    // 传入非负整数 k，返回当前方阵的 k 次幂；零次幂为单位矩阵。
    // 要求当前矩阵为方阵，时间为 $O(n^3\log k)$。
    Matrix pow(ll k) const {
        assert(n == m && k >= 0);
        Matrix ret = identity(n), mul = *this;
        for (; k; k >>= 1, mul = mul * mul) {
            if (k & 1) ret = ret * mul;
        }
        return ret;
    }
};
\end{lstlisting}

\subsection{异或线性基（参考）}
\noindent{\small \texttt{Math/Alg/\_linear\_basis.hpp}}\par

\begin{lstlisting}
#pragma once
#include "../head.hpp"

// 维护 B 位无符号整数的子集异或值；B 的范围为 $1\le B\le64$。
// 默认构造空线性基；p 存基向量，rank 是秩，zero 表示能否由非空子集异或得到零。
template <int B = 64>
struct LinearBasis {
    static_assert(1 <= B && B <= 64);
    array<u64, B> p{};
    int rank = 0;
    bool zero = false; 

    // 传入一个 B 位整数 x，将它并入线性基；返回是否增加了秩。
    // 若 x 与已有元素线性相关，则记录存在非空零异或子集；时间为 $O(B)$。
    bool insert(u64 x) {
        if constexpr (B < 64) assert((x >> B) == 0);
        for (int i = B - 1; i >= 0; i--) {
            if (!(x >> i & 1)) continue;
            if (p[i]) x ^= p[i];
            else {
                p[i] = x;
                rank++;
                return true;
            }
        }
        zero = true;
        return false;
    }
    // 传入 x，返回它能否由当前元素的某个子集异或得到，允许空集。
    bool contains(u64 x) const {
        for (int i = B - 1; i >= 0; i--) {
            if (x >> i & 1) x ^= p[i];
        }
        return x == 0;
    }
    // 传入初始值 x，返回 x 与任意子集异或后的最大值；默认从零开始。
    u64 max_xor(u64 x = 0) const {
        for (int i = B - 1; i >= 0; i--) x = max(x, x ^ p[i]);
        return x;
    }
    // 传入初始值 x，返回 x 与任意子集异或后的最小值；默认从零开始。
    u64 min_xor(u64 x = 0) const {
        for (int i = B - 1; i >= 0; i--) x = min(x, x ^ p[i]);
        return x;
    }
    // 传入另一个线性基 b，将其代表的元素并入当前对象，不返回新对象。
    void merge(const LinearBasis& b) {
        auto q = b.p;
        zero |= b.zero;
        for (u64 x : q) if (x) insert(x);
    }
    
    // 无需参数，返回按主元从低到高排列、已消去低主元的基向量。
    // 该数组按二进制编号组合时，异或值递增；构造时间为 $O(B^2)$，可缓存用于多次查询。
    vector<u64> ordered() const {
        auto q = p;
        vector<u64> ret;
        for (int i = 0; i < B; i++) {
            if (!q[i]) continue;
            for (int j = 0; j < i; j++) {
                if (q[i] >> j & 1) q[i] ^= q[j];
            }
            ret.push_back(q[i]);
        }
        return ret;
    }

    // 传入从零开始的编号 k，返回第 k 小的不同异或值；越界返回 nullopt。
    // nonempty 为真时排除空集；若非空子集也能得到零，仍保留零。单次时间为 $O(B^2)$。
    optional<u64> kth(u64 k, bool nonempty = false) const {
        if (nonempty && !zero) {
            if (k == numeric_limits<u64>::max()) return nullopt;
            k++;
        }
        if (rank < 64 && k >= (u64(1) << rank)) return nullopt;
        auto q = ordered();
        u64 ret = 0;
        for (int i = 0; i < rank; i++) if (k >> i & 1) ret ^= q[i];
        return ret;
    }
};
\end{lstlisting}

\subsection{行列式（参考）}
\noindent{\small \texttt{Math/Alg/\_determinant.hpp}}\par

\begin{lstlisting}
#pragma once
#include "../head.hpp"
#include "../NT/modint.hpp"

// 传入方阵 a，返回其行列式；空矩阵返回一，不修改传入的原数组。
// Z 必须是域元素，通常使用质数模 mint；不适用于整数除法或一般合数模。时间为 $O(n^3)$。
template <typename Z>
Z determinant(vector<vector<Z>> a) {
    int n = a.size();
    for (const auto& row : a) assert((int)row.size() == n);
    Z ret = 1;
    for (int col = 0; col < n; col++) {
        int pivot = col;
        while (pivot < n && a[pivot][col] == Z(0)) pivot++;
        if (pivot == n) return 0;
        if (pivot != col) {
            swap(a[pivot], a[col]);
            ret = -ret;
        }
        ret *= a[col][col];
        Z iv = Z(1) / a[col][col];
        for (int i = col + 1; i < n; i++) {
            Z t = a[i][col] * iv;
            for (int j = col + 1; j < n; j++) a[i][j] -= t * a[col][j];
            a[i][col] = 0;
        }
    }
    return ret;
}
\end{lstlisting}

\subsection{矩阵求逆（参考）}
\noindent{\small \texttt{Math/Alg/\_matrix\_inverse.hpp}}\par

\begin{lstlisting}
#pragma once
#include "../head.hpp"
#include "../NT/modint.hpp"

// 传入方阵 a，返回 {是否可逆, 逆矩阵}；不可逆时返回 {false, {}}。
// 空矩阵返回 {true, {}}。Z 须为域元素，通常使用质数模 mint；时间为 $O(n^3)$。
template <typename Z>
pair<bool, vector<vector<Z>>> matrix_inv(vector<vector<Z>> a) {
    int n = a.size();
    for (int i = 0; i < n; i++) {
        assert((int)a[i].size() == n);
        a[i].resize(2 * n);
        a[i][n + i] = 1;
    }
    for (int col = 0; col < n; col++) {
        int pivot = col;
        while (pivot < n && a[pivot][col] == Z(0)) pivot++;
        if (pivot == n) return {false, {}};
        swap(a[pivot], a[col]);
        Z iv = Z(1) / a[col][col];
        for (int j = col; j < 2 * n; j++) a[col][j] *= iv;
        for (int i = 0; i < n; i++) {
            if (i == col || a[i][col] == Z(0)) continue;
            Z t = a[i][col];
            for (int j = col; j < 2 * n; j++) a[i][j] -= t * a[col][j];
        }
    }
    for (auto& row : a) row.erase(row.begin(), row.begin() + n);
    return {true, a};
}
\end{lstlisting}

\subsection{异或高斯消元（参考）}
\noindent{\small \texttt{Math/Alg/\_gauss\_xor.hpp}}\par

\begin{lstlisting}
#pragma once
#include "../head.hpp"

// 求解异或方程组：a 是增广矩阵，m 是未知数数量，要求 $0\le m<M$。
// 前 m 位为系数，第 m 位为常数；返回 {rank, 一个特解}，自由变量取零，无解时 rank 为 -1。
// 每次整行消元使用 bitset；时间约为 $O(n\min(n,m)M/w)$，w 为机器字长。
template <size_t M>
pair<int, vector<int>> gauss_xor(vector<bitset<M>> a, int m) {
    assert(0 <= m && m < (int)M);
    int n = a.size(), rank = 0;
    vector<int> pos;
    for (int col = 0; col < m && rank < n; col++) {
        int pivot = rank;
        while (pivot < n && !a[pivot][col]) pivot++;
        if (pivot == n) continue;
        swap(a[pivot], a[rank]);
        for (int i = rank + 1; i < n; i++) {
            if (a[i][col]) a[i] ^= a[rank];
        }
        pos.push_back(col);
        rank++;
    }
    for (int i = rank; i < n; i++) {
        if (a[i][m]) return {-1, {}};
    }
    bitset<M> x;
    for (int i = rank - 1; i >= 0; i--) {
        x[pos[i]] = a[i][m] ^ ((a[i] & x).count() & 1);
    }
    vector<int> ret(m);
    for (int i = 0; i < m; i++) ret[i] = x[i];
    return {rank, ret};
}
\end{lstlisting}

\subsection{矩阵树定理（参考）}
\noindent{\small \texttt{Math/Alg/\_kirchhoff.hpp}}\par

\begin{lstlisting}
#pragma once
#include "_determinant.hpp"

// 传入点数 n 和无向边 edges，返回所有生成树的边权乘积之和。
// 点从零编号，n 至少为一，边为 {u, v, w}；允许重边，自环忽略，普通计数令边权为一。
// Z 通常取质数模 mint；单点图返回一，时间为 $O(n^3)$。
template <typename Z>
Z spanning_trees(int n, const vector<tuple<int, int, Z>>& edges) {
    assert(n >= 1);
    vector<vector<Z>> a(n - 1, vector<Z>(n - 1));
    for (auto [u, v, w] : edges) {
        if (u == v) continue;
        if (u < n - 1) a[u][u] += w;
        if (v < n - 1) a[v][v] += w;
        if (u < n - 1 && v < n - 1) {
            a[u][v] -= w;
            a[v][u] -= w;
        }
    }
    return determinant(move(a));
}

// 传入点数 n、有向边 edges 和根 root，返回所有朝向根的生成树的权值和。
// 边 {u, v, w} 表示从 u 到 v；每个点须沿树边最终到达 root。点从零编号，Z 为域元素。
template <typename Z>
Z arborescences(int n, const vector<tuple<int, int, Z>>& edges, int root) {
    assert(0 <= root && root < n);
    vector<vector<Z>> a(n - 1, vector<Z>(n - 1));
    // 传入原点编号 x，返回删除 root 后对应的矩阵下标；调用时 x 不能是 root。
    auto id = [&](int x) { return x - (x > root); };
    for (auto [u, v, w] : edges) {
        if (u == v || u == root) continue;
        a[id(u)][id(u)] += w;
        if (v != root) a[id(u)][id(v)] -= w;
    }
    return determinant(move(a));
}
\end{lstlisting}

\subsection{自适应 Simpson 积分（参考）}
\noindent{\small \texttt{Math/Alg/\_integrate.hpp}}\par

\begin{lstlisting}
#pragma once
#include "../head.hpp"

// 传入函数 f、积分端点 l 和 r，返回定积分的 long double 估计值。
// eps 是目标绝对误差，depth 是递归深度上限；f 应在闭区间内有限且足够光滑。
// 不连续点需分段，端点奇点需换元；到达深度上限时返回当前估计，不保证已满足 eps。
template <typename F>
long double integrate(F f, long double l, long double r,
                      long double eps = 1e-9L, int depth = 25) {
    using Real = long double;
    assert(eps > 0 && depth >= 0);
    if (l == r) return 0;
    if (l > r) return -integrate(f, r, l, eps, depth);
    // 传入区间端点 l、r，以及左端点、中点、右端点的函数值 a、b、c，返回 Simpson 估计。
    auto simpson = [](Real l, Real r, Real a, Real b, Real c) {
        return (r - l) * (a + 4 * b + c) / 6;
    };
    // 接收区间、已缓存的函数值 a、b、c、旧估计 val、误差额度 eps 和剩余深度 depth。
    // 返回细分后的积分估计；self 用于递归，两侧区间分别使用一半误差额度。
    auto work = [&](auto self, Real l, Real r, Real a, Real b, Real c,
                    Real val, Real eps, int depth) -> Real {
        Real mid = l + (r - l) / 2;
        Real lm = l + (mid - l) / 2, rm = mid + (r - mid) / 2;
        if (lm == l || rm == r) return val;
        Real d = f(lm), e = f(rm);
        Real lt = simpson(l, mid, a, d, b), rt = simpson(mid, r, b, e, c);
        Real delta = lt + rt - val;
        if (depth == 0 || abs(delta) <= 15 * eps) return lt + rt + delta / 15;
        return self(self, l, mid, a, d, b, lt, eps / 2, depth - 1)
             + self(self, mid, r, b, e, c, rt, eps / 2, depth - 1);
    };
    Real mid = l + (r - l) / 2, a = f(l), b = f(mid), c = f(r);
    return work(work, l, r, a, b, c, simpson(l, r, a, b, c), eps, depth);
}
\end{lstlisting}

\subsection{Nim 与 SG（参考）}
\noindent{\small \texttt{Math/Alg/\_sg.hpp}}\par

\begin{lstlisting}
#pragma once
#include "../head.hpp"

// 传入各堆石子数 a，判断普通 Nim 是否先手必胜；空局面返回 false。
inline bool nim_win(const vector<u64>& a) {
    u64 ret = 0;
    for (u64 x : a) ret ^= x;
    return ret != 0;
}
// 传入各堆石子数 a，判断拿走最后一颗者输的反常 Nim 是否先手必胜。
// 无石子局面按反常规则返回 true；不能将此规则直接用于一般 SG 游戏。
inline bool misere_nim_win(const vector<u64>& a) {
    u64 ret = 0;
    int ones = 0;
    bool big = false;
    for (u64 x : a) {
        ret ^= x;
        ones += x == 1;
        big |= x > 1;
    }
    return big ? ret != 0 : ones % 2 == 0;
}

// 传入无环局面图 g，g[u] 列出局面 u 一步可到达的局面；返回各点的 SG 值。
// 点从零编号，不能行动者输；若存在环则断言失败。时间为 $O(n+m)$。
inline vector<int> sg_dag(const vector<vector<int>>& g) {
    int n = g.size(), cnt = 0;
    vector<vector<int>> rev(n);
    vector<int> deg(n), sg(n), seen;
    queue<int> q;
    for (int u = 0; u < n; u++) {
        deg[u] = g[u].size();
        for (int v : g[u]) rev[v].push_back(u);
        if (deg[u] == 0) q.push(u);
    }
    while (!q.empty()) {
        int u = q.front();
        q.pop();
        cnt++;
        seen.assign(g[u].size() + 1, 0);
        for (int v : g[u]) {
            if (sg[v] < (int)seen.size()) seen[sg[v]] = 1;
        }
        while (seen[sg[u]]) sg[u]++;
        for (int v : rev[u]) if (--deg[v] == 0) q.push(v);
    }
    assert(cnt == n);
    return sg;
}
\end{lstlisting}


\providecommand{\tightlist}{\setlength{\itemsep}{0pt}\setlength{\parskip}{0pt}}
\providecommand{\passthrough}[1]{#1}
\hypertarget{ux6570ux5b66ux77e5ux8bc6ux70b9}{%
\subsection{数学知识点}\label{ux6570ux5b66ux77e5ux8bc6ux70b9}}

本文整理赛时常用的数学结论、推导方法和使用条件，既包含模板的配套说明，也包含无需单独写成模板的知识点。函数的参数与返回值见各头文件的注释，调用示例见 \href{_templates.md}{模板索引}。下划线文件仍保留待审核标记。

\hypertarget{ux7ec4ux5408ux8ba1ux6570}{%
\subsubsection{1. 组合计数}\label{ux7ec4ux5408ux8ba1ux6570}}

\hypertarget{ux7ec4ux5408ux6570ux591aux91cdux96c6ux6392ux5217ux4e0eux9694ux677fux6cd5}{%
\paragraph{组合数、多重集排列与隔板法}\label{ux7ec4ux5408ux6570ux591aux91cdux96c6ux6392ux5217ux4e0eux9694ux677fux6cd5}}

从 \(n\) 个不同元素中选出 \(r\) 个，方案数为 \(\binom{n}{r}\)。当 \(r<0\) 或 \(r>n\) 时，约定组合数为 \(0\)。

若第 \(i\) 类元素有 \(c_i\) 个，同类元素不区分，则排列数为

\[
\frac{\left(\sum_{i=1}^{k}c_i\right)!}{\prod_{i=1}^{k}c_i!}.
\]

对于非负整数方程

\[
x_1+\cdots+x_k=n,
\]

当 \(k\ge1\) 时，解数为 \(\binom{n+k-1}{k-1}\)。当 \(k=0\) 时，只有 \(n=0\) 才有一个解。

若要求 \(x_i\ge l_i\)，先令 \(y_i=x_i-l_i\)，再使用隔板法。若还要求 \(x_i\le u_i\)，可以对违反上界的变量容斥；选中变量 \(i\) 时，将剩余总量减去 \(u_i+1\)。同时有下界时，应先平移变量，再使用平移后的上界。

使用阶乘与逆阶乘计算组合数时，通常要求模数 \(P\) 为质数，且预处理上限小于 \(P\)。Lucas 定理适用于很大的 \(n\)，但每一位组合数仍需实际计算；模数很大时，不能只看递归层数。合数模可使用 exLucas，当前版本的时间和空间取决于各素数幂的大小。

\hypertarget{ux7ec4ux5408ux6052ux7b49ux5f0fux4e0eux4e8cux9879ux5f0fux53cdux6f14}{%
\paragraph{组合恒等式与二项式反演}\label{ux7ec4ux5408ux6052ux7b49ux5f0fux4e0eux4e8cux9879ux5f0fux53cdux6f14}}

下列公式中，上标为非负整数，下标越界的组合数取 \(0\)：

\[
\binom{n}{k}=\binom{n}{n-k},\qquad
\binom{n}{k}=\binom{n-1}{k-1}+\binom{n-1}{k}.
\]

第二个等式要求 \(n\ge1\)。常用求和形式为

\[
\sum_{k=0}^{n}\binom{n}{k}=2^n,\qquad
\sum_{k=0}^{n}(-1)^k\binom{n}{k}=[n=0],
\]

\[
\sum_{i=r}^{n}\binom{i}{r}=\binom{n+1}{r+1}\quad(0\le r\le n),
\]

\[
\sum_{k=0}^{r}\binom{n}{k}\binom{m}{r-k}=\binom{n+m}{r}.
\]

最后一个等式是 Vandermonde 卷积；取 \(m=n,r=n\)，得到 \(\sum_{k=0}^{n}\binom{n}{k}^2=\binom{2n}{n}\)。遇到带 \(k\) 的系数，可先用 \(k\binom{n}{k}=n\binom{n-1}{k-1}\) 消掉它，其中 \(n\ge1\)。

二项式反演为

\[
b_n=\sum_{k=0}^{n}\binom{n}{k}a_k
\quad\Longleftrightarrow\quad
 a_n=\sum_{k=0}^{n}(-1)^{n-k}\binom{n}{k}b_k.
\]

它常用于从``选出一个子集后的总计数''还原原始计数。模意义下，这些不含除法的整数恒等式对任意模数都成立；计算组合数的方法仍需满足各自条件。

\hypertarget{ux5bb9ux65a5}{%
\paragraph{容斥}\label{ux5bb9ux65a5}}

有限集合的并集满足

\[
\left|\bigcup_{i=1}^{k}A_i\right|
=\sum_{\varnothing\ne S\subseteq\{1,\ldots,k\}}
(-1)^{|S|+1}\left|\bigcap_{i\in S}A_i\right|.
\]

将 \(n\) 个不同的球放入 \(k\) 个不同的盒子，并要求每个盒子非空，方案数为

\[
\sum_{i=0}^{k}(-1)^i\binom{k}{i}(k-i)^n=k!S(n,k).
\]

这里 \(S(n,k)\) 是第二类斯特林数。盒子有标号，所以需要乘上 \(k!\)。

若有 \(k\) 个性质，令 \(F_j\) 为任取 \(j\) 个性质时，满足它们的对象数量之和，则恰好满足 \(r\) 个性质的对象数为

\[
\sum_{j=r}^{k}(-1)^{j-r}\binom{j}{r}F_j.
\]

在这类计数公式中通常约定 \(0^0=1\)。例如，空集到空集恰好有一个映射。

\hypertarget{ux53cdux5c04ux539fux7406ballot-ux6570ux4e0e-catalan-ux6570}{%
\paragraph{反射原理、Ballot 数与 Catalan 数}\label{ux53cdux5c04ux539fux7406ballot-ux6570ux4e0e-catalan-ux6570}}

考虑由 \(a\) 次 \(+1\) 和 \(b\) 次 \(-1\) 组成的序列，要求任意前缀和非负。当 \(a<b\) 时无解；当 \(a\ge b\) 时，方案数为

\[
\binom{a+b}{b}-\binom{a+b}{b-1}.
\]

证明时，将不合法路径第一次到达 \(-1\) 之前的所有步反向，就能把不合法路径与另一类容易计数的路径一一对应。这是反射原理的常见用法。

若要求每个非空前缀和都严格为正，且 \(a>b\)，则方案数为

\[
\frac{a-b}{a+b}\binom{a+b}{b}.
\]

模意义下不一定能直接除以 \(a+b\)。也可以先固定首步为 \(+1\)，然后把其余部分转化为前缀和非负的问题。

Catalan 数满足

\[
C_n=\binom{2n}{n}-\binom{2n}{n-1},\qquad C_0=1.
\]

它可以计数合法括号串、不越过对角线的格点路径、二叉树形状和凸多边形三角剖分。差分形式避免了单独除以 \(n+1\)，但所用的组合数函数仍需满足预处理范围等条件。

\hypertarget{ux9519ux6392ux65afux7279ux6797ux6570ux4e0e-bell-ux6570}{%
\paragraph{错排、斯特林数与 Bell 数}\label{ux9519ux6392ux65afux7279ux6797ux6570ux4e0e-bell-ux6570}}

错排指没有任何元素停留在原位置的排列。其数量满足

\[
D_0=1,\qquad D_1=0,\qquad
D_n=(n-1)(D_{n-1}+D_{n-2})\quad(n\ge2).
\]

无符号第一类斯特林数 \(c(n,k)\) 统计恰有 \(k\) 个置换环的 \(n\) 元排列：

\[
c(n,k)=c(n-1,k-1)+(n-1)c(n-1,k).
\]

第二类斯特林数 \(S(n,k)\) 统计将 \(n\) 个不同元素划分成 \(k\) 个非空、无标号集合的方案：

\[
S(n,k)=S(n-1,k-1)+kS(n-1,k).
\]

两者都以 \(c(0,0)=S(0,0)=1\) 为初值；当 \(n>0\) 时，第 \(0\) 列为 \(0\)。

Bell 数允许集合数量任意，因此

\[
B_n=\sum_{k=0}^{n}S(n,k),\qquad B_0=1.
\]

若需要快速计算第二类斯特林数的一行，可使用

\[
S(n,k)=\sum_{i=0}^{k}\frac{i^n}{i!}\frac{(-1)^{k-i}}{(k-i)!}.
\]

右侧是一次卷积。计算复杂度时，还需计入各个 \(i^n\) 的快速幂开销；阶乘的逆元也必须存在。

\hypertarget{ux6574ux6570ux62c6ux5206ux4e0e-euler-ux4e94ux8fb9ux5f62ux6570ux5b9aux7406}{%
\paragraph{整数拆分与 Euler 五边形数定理}\label{ux6574ux6570ux62c6ux5206ux4e0e-euler-ux4e94ux8fb9ux5f62ux6570ux5b9aux7406}}

普通五边形数为 \(j(3j-1)/2\)，其中 \(j\ge0\)。将 \(j\) 扩展为所有整数，得到广义五边形数；正数部分按大小排列为

\[
1,2,5,7,12,15,22,26,\ldots.
\]

实现时通常成对枚举 \(j(3j-1)/2\) 和 \(j(3j+1)/2\)，其中 \(j\ge1\)。Euler 五边形数定理给出形式幂级数恒等式

\[
\prod_{k\ge1}(1-x^k)
=\sum_{j\in\mathbb Z}(-1)^j x^{j(3j-1)/2}
=1+\sum_{j\ge1}(-1)^j
 \left(x^{j(3j-1)/2}+x^{j(3j+1)/2}\right).
\]

令 \(p(n)\) 表示将 \(n\) 拆成若干个正整数之和的方案数，不区分部件顺序。约定 \(p(0)=1\)，负下标的值为 \(0\)。其生成函数为

\[
P(x)=\sum_{n\ge0}p(n)x^n=\prod_{k\ge1}\frac{1}{1-x^k}.
\]

两式相乘后比较系数，得到五边形数递推

\[
p(n)=\sum_{j\ge1}(-1)^{j-1}
\left[
 p\left(n-\frac{j(3j-1)}{2}\right)
 +p\left(n-\frac{j(3j+1)}{2}\right)
\right].
\]

递推中的符号依次为 \(+,+,-,-,+,+,-,-,\ldots\)，只枚举不超过 \(n\) 的广义五边形数。计算前 \(n\) 项需要 \(O(n\sqrt n)\) 时间和 \(O(n)\) 空间，对任意正整数模数都成立；实现见 \passthrough{\lstinline!Comb/\_numbers.hpp!} 的 \passthrough{\lstinline!partition\_numbers!}。

常见的拆分限制可以直接写成生成函数：

\begin{itemize}
\tightlist
\item
  每个大小 \(k\) 可重复使用：因子为 \((1-x^k)^{-1}\)。
\item
  每个大小 \(k\) 最多用一次：因子为 \(1+x^k\)。
\item
  每个大小 \(k\) 最多用 \(c_k\) 次：因子为 \(1+x^k+\cdots+x^{c_k k}\)。
\item
  恰好拆成 \(k\) 个正整数：从每个部件减一，等价于把 \(n-k\) 拆成至多 \(k\) 个正整数。
\item
  至多拆成 \(k\) 个正整数的方案数，等于每个部件不超过 \(k\) 的方案数，可由 Ferrers 图转置得到。
\end{itemize}

不同部件拆分与奇数部件拆分数量相同，因为

\[
\prod_{k\ge1}(1+x^k)
=\prod_{k\ge1}\frac{1-x^{2k}}{1-x^k}
=\prod_{k\ge1}\frac{1}{1-x^{2k-1}}.
\]

上述讨论都不区分部件顺序。若区分顺序，正整数 \(n\) 的无约束拆分数为 \(2^{n-1}\)；空拆分的数量为 \(1\)。只处理有限部件或次数限制时，背包通常更容易实现。

\hypertarget{burnsideux9879ux94feux4e0eux624bux94fe}{%
\subsubsection{2. Burnside、项链与手链}\label{burnsideux9879ux94feux4e0eux624bux94fe}}

先明确哪些变换会让两个方案被视为相同。这些变换组成群 \(G\)，它们将方案集合 \(X\) 分成若干个等价类，也称轨道。Burnside 引理给出轨道数：

\[
|X/G|=\frac{1}{|G|}\sum_{g\in G}|\operatorname{Fix}(g)|.
\]

其中 \(\operatorname{Fix}(g)\) 表示经过变换 \(g\) 后保持不变的方案。使用公式时，需要对完整的群求和，不能只对几个生成元取平均。

若位置置换有 \(c\) 个环，每个位置可以独立使用 \(q\) 种颜色，则不动染色数为 \(q^c\)。这是因为同一置换环内的颜色必须相同。若题目还限制相邻位置的颜色，就需要重新计算不动方案数。

\hypertarget{ux9879ux94feux53eaux5c06ux65cbux8f6cux89c6ux4e3aux76f8ux540c}{%
\paragraph{项链：只将旋转视为相同}\label{ux9879ux94feux53eaux5c06ux65cbux8f6cux89c6ux4e3aux76f8ux540c}}

对于 \(n\ge1\) 个位置和 \(q\) 种颜色，项链数量为

\[
N(n,q)=\frac{1}{n}\sum_{d\mid n}\varphi(d)q^{n/d}.
\]

旋转 \(k\) 格会产生 \(\gcd(n,k)\) 个置换环，按最大公约数分类即可得到该式。

\hypertarget{ux624bux94feux540cux65f6ux5c06ux65cbux8f6cux4e0eux7ffbux8f6cux89c6ux4e3aux76f8ux540c}{%
\paragraph{手链：同时将旋转与翻转视为相同}\label{ux624bux94feux540cux65f6ux5c06ux65cbux8f6cux4e0eux7ffbux8f6cux89c6ux4e3aux76f8ux540c}}

令 \(R(n,q)\) 表示所有旋转的不动方案数之和，所有翻转的不动方案数之和为

\[
F(n,q)=
\begin{cases}
 n q^{(n+1)/2}, & n\text{ 为奇数},\\
 \dfrac{n}{2}\left(q^{n/2+1}+q^{n/2}\right), & n\text{ 为偶数}.
\end{cases}
\]

手链数量为

\[
H(n,q)=\frac{R(n,q)+F(n,q)}{2n}.
\]

这些平均公式在 \(n=1\) 和 \(n=2\) 时仍成立，尽管某些不同的群元素会诱导相同的位置置换。

\hypertarget{ux6a21ux610fux4e49ux4e0bux7684ux9664ux6cd5}{%
\paragraph{模意义下的除法}\label{ux6a21ux610fux4e49ux4e0bux7684ux9664ux6cd5}}

只有当分母可逆时，才能直接乘逆元。若质数模 \(P\) 整除 \(|G|\)，不能直接使用 \passthrough{\lstinline!Z(group.size())!} 的逆元。

若能在模数 \(|G|P\) 下准确求出分子的余数，可以先对该余数做整数除法，再得到答案模 \(P\) 的结果。这里需要另外处理扩大后的模数和乘法范围。

当前 \passthrough{\lstinline!necklaces!} 和 \passthrough{\lstinline!bracelets!} 使用试除分解 \(n\)，这部分最坏需要 \(O(\sqrt n)\) 时间。若 \(n\) 很大，可将分解方法改为 \passthrough{\lstinline!rho::prime\_factor!}。

\hypertarget{ux6807ux53f7ux6811pruxfcfer-ux5e8fux5217ux4e0eux77e9ux9635ux6811ux5b9aux7406}{%
\subsubsection{3. 标号树、Prüfer 序列与矩阵树定理}\label{ux6807ux53f7ux6811pruxfcfer-ux5e8fux5217ux4e0eux77e9ux9635ux6811ux5b9aux7406}}

\hypertarget{pruxfcfer-ux5e8fux5217ux4e0eux6811ux8ba1ux6570}{%
\paragraph{Prüfer 序列与树计数}\label{pruxfcfer-ux5e8fux5217ux4e0eux6811ux8ba1ux6570}}

当 \(n\ge2\) 时，标号树与长度为 \(n-2\)、元素取自 \(\{0,\ldots,n-1\}\) 的 Prüfer 序列一一对应。编码时，每次删除编号最小的叶子，并记录它的邻点。点 \(i\) 在序列中恰好出现 \(d_i-1\) 次，其中 \(d_i\) 是它的度数。

由此得到以下结论：

\begin{itemize}
\tightlist
\item
  \(n\) 个标号点的无根树数为 \(n^{n-2}\)；单点树的数量单独取 \(1\)。
\item
  给定度数 \(d_i\ge1\)，且 \(\sum_{i=1}^{n}d_i=2n-2\)，树数为 \(\dfrac{(n-2)!}{\prod_{i=1}^{n}(d_i-1)!}\)。
\item
  单点树的度数为 \(0\)。空 Prüfer 序列在当前解码接口中表示两个点的树。
\item
  若固定森林有 \(k\ge2\) 个连通块，大小分别为 \(s_1,\ldots,s_k\)，则包含该森林的标号生成树数为 \(n^{k-2}\prod_{i=1}^{k}s_i\)。当 \(k=1\) 时，数量为 \(1\)。
\end{itemize}

\hypertarget{ux65e0ux5411ux77e9ux9635ux6811ux5b9aux7406}{%
\paragraph{无向矩阵树定理}\label{ux65e0ux5411ux77e9ux9635ux6811ux5b9aux7406}}

构造 Laplace 矩阵 \(L\)：对角线存放各点的加权度数，非对角线 \(L_{u,v}\) 存放点 \(u,v\) 之间所有边权之和的相反数。删除同一编号的一行和一列，所得行列式就是生成树权值之和；每棵树的权值是其边权之积。

允许重边，自环不影响答案。单点图的答案为 \(1\)，对应空矩阵行列式为 \(1\)。

\hypertarget{ux6709ux5411ux77e9ux9635ux6811ux5b9aux7406}{%
\paragraph{有向矩阵树定理}\label{ux6709ux5411ux77e9ux9635ux6811ux5b9aux7406}}

当前 \passthrough{\lstinline!arborescences!} 将出边权值之和放在对角线，并对每条边 \(u\to v\) 从 \(L_{u,v}\) 中减去边权。删除根对应的行和列后，统计的是每个点沿树边最终到达根的树。如果需要统计从根能够到达所有点的树，将所有边反向后调用。

质数模下得到的是计数的模值。行列式模 \(P\) 为 \(0\)，并不能说明原图一定不连通。当前除法消元也不适用于一般合数模。

\hypertarget{ux57faux7840ux6570ux8bbaux83abux6bd4ux4e4cux65afux53cdux6f14ux4e0eux6570ux8bbaux5206ux5757}{%
\subsubsection{4. 基础数论、莫比乌斯反演与数论分块}\label{ux57faux7840ux6570ux8bbaux83abux6bd4ux4e4cux65afux53cdux6f14ux4e0eux6570ux8bbaux5206ux5757}}

\hypertarget{ux8d28ux56e0ux6570ux5206ux89e3ux7ea6ux6570ux4e0eux9636ux4e58ux4e2dux7684ux7d20ux56e0ux5b50}{%
\paragraph{质因数分解、约数与阶乘中的素因子}\label{ux8d28ux56e0ux6570ux5206ux89e3ux7ea6ux6570ux4e0eux9636ux4e58ux4e2dux7684ux7d20ux56e0ux5b50}}

若 \(n=\prod_{i=1}^{r}p_i^{e_i}\ge1\)，则约数个数与约数和分别为

\[
\tau(n)=\prod_{i=1}^{r}(e_i+1),\qquad
\sigma(n)=\prod_{i=1}^{r}(1+p_i+\cdots+p_i^{e_i}).
\]

约数和中的每个因子也可写成 \((p_i^{e_i+1}-1)/(p_i-1)\)，但模意义下分母可能不可逆。没有逆元时，应直接求几何级数。

记 \(v_p(a)\) 为正整数 \(a\) 中素因子 \(p\) 的指数。Legendre 公式为

\[
v_p(n!)=\sum_{j\ge1}\left\lfloor\frac{n}{p^j}\right\rfloor,
\]

因此

\[
v_p\!\left(\binom{n}{k}\right)
=v_p(n!)-v_p(k!)-v_p((n-k)!).
\]

循环中不断令 \(n\gets\lfloor n/p\rfloor\) 再累加，可以避免显式计算 \(p^j\) 溢出。

\hypertarget{ux7ebfux6027ux540cux4f59ux4e0eux9ea6ux4e50ux9e21ux5b9aux7406}{%
\paragraph{线性同余与麦乐鸡定理}\label{ux7ebfux6027ux540cux4f59ux4e0eux9ea6ux4e50ux9e21ux5b9aux7406}}

对于 \(m\ge1\)，同余方程 \(ax\equiv b\pmod m\) 有解，当且仅当 \(g=\gcd(a,m)\) 整除 \(b\)。有解时，先将 \(a,b,m\) 同除以 \(g\)，再求逆元；模 \(m\) 的不同解恰有 \(g\) 个。

若互质整数 \(a,b\ge2\)，则不能表示为 \(ax+by\)、其中 \(x,y\ge0\) 的最大非负整数为

\[
ab-a-b.
\]

不能表示的非负整数共有 \((a-1)(b-1)/2\) 个。若 \(a\) 或 \(b\) 为 \(1\)，所有非负整数都能表示；若二者不互质，不可表示的非负整数有无穷多个。不要将两种面额的公式直接用于三种及以上面额。

\hypertarget{ux6b27ux62c9ux51fdux6570ux9636ux4e0eux539fux6839}{%
\paragraph{欧拉函数、阶与原根}\label{ux6b27ux62c9ux51fdux6570ux9636ux4e0eux539fux6839}}

对于 \(n\ge1\)，有

\[
\varphi(n)=n\prod_{p\mid n}\left(1-\frac1p\right),\qquad
\varphi(1)=1.
\]

当 \(\gcd(a,b)=1\) 时，\(\varphi(ab)=\varphi(a)\varphi(b)\)。若素数 \(p\mid n\)，则

\[
\varphi(n)=
\begin{cases}
 p\varphi(n/p), & p^2\mid n,\\
 (p-1)\varphi(n/p), & p^2\nmid n.
\end{cases}
\]

当 \(n\ge2\) 时，\([1,n]\) 中与 \(n\) 互质的整数之和为 \(n\varphi(n)/2\)；\(n=1\) 时应单独取 \(1\)。另有

\[
\sum_{i=1}^{n}\gcd(i,n)=\sum_{d\mid n}d\varphi(n/d).
\]

若 \(m\ge2\) 且 \(\gcd(a,m)=1\)，Euler 定理给出 \(a^{\varphi(m)}\equiv1\pmod m\)。\(a\) 模 \(m\) 的阶是使 \(a^d\equiv1\pmod m\) 的最小正整数 \(d\)，且 \(d\mid\varphi(m)\)。因此互质时可以将指数对阶取模；不互质时不能直接这样处理。

对于质数 \(p\)，原根 \(g\) 的阶为 \(p-1\)。分解 \(p-1\) 后，对其每个不同素因子 \(q\) 检查

\[
g^{(p-1)/q}\not\equiv1\pmod p
\]

即可判定非零剩余类 \(g\) 是否为原根。阶为 \(N\) 的 NTT 单位根要求 \(N\mid p-1\)，再由原根构造 \(g^{(p-1)/N}\)。

\hypertarget{ux83abux6bd4ux4e4cux65afux53cdux6f14ux4e0eux5e38ux7528ux6052ux7b49ux5f0f}{%
\paragraph{莫比乌斯反演与常用恒等式}\label{ux83abux6bd4ux4e4cux65afux53cdux6f14ux4e0eux5e38ux7528ux6052ux7b49ux5f0f}}

对于正整数 \(n\)，有

\[
\sum_{d\mid n}\mu(d)=[n=1],\qquad
\sum_{d\mid n}\varphi(d)=n.
\]

方括号 \([n=1]\) 在条件成立时取 \(1\)，否则取 \(0\)。

约数形式的莫比乌斯反演为

\[
F(n)=\sum_{d\mid n}f(d)
\quad\Longleftrightarrow\quad
f(n)=\sum_{d\mid n}\mu(d)F(n/d).
\]

对于只在有限范围内非零的函数，倍数形式为

\[
F(n)=\sum_{n\mid d}f(d)
\quad\Longleftrightarrow\quad
f(n)=\sum_{n\mid d}\mu(d/n)F(d).
\]

当 \(n,m\ge0\) 时，来自 \([1,n]\times[1,m]\) 的互质有序数对数量为

\[
\sum_{d=1}^{\min(n,m)}\mu(d)
\left\lfloor\frac{n}{d}\right\rfloor
\left\lfloor\frac{m}{d}\right\rfloor.
\]

这些数对的最大公约数之和为

\[
\sum_{d=1}^{\min(n,m)}\varphi(d)
\left\lfloor\frac{n}{d}\right\rfloor
\left\lfloor\frac{m}{d}\right\rfloor.
\]

若题目只统计无序数对，或要求 \(i<j\)，还需处理对称性和对角线。

\hypertarget{ux5546ux5206ux5757}{%
\paragraph{商分块}\label{ux5546ux5206ux5757}}

对于 \(l\ge1\)，设

\[
q=\left\lfloor\frac{n}{l}\right\rfloor,\qquad
r=\left\lfloor\frac{n}{q}\right\rfloor.
\]

区间 \([l,r]\) 内的商相同。处理完这一段后，从 \(r+1\) 继续，总共只有 \(O(\sqrt n)\) 段。

若同时处理 \(\lfloor n/l\rfloor\) 与 \(\lfloor m/l\rfloor\)，右端点取两个分块右端点的较小值。若权函数有前缀和，就可以一次求出整段权重，再乘上两个商。

当 \(n=0\) 时没有分块。起点不能设为 \(0\)。若参数可能达到 \passthrough{\lstinline!ll!} 的最大值，还需避免最后一段结束后计算溢出的 \passthrough{\lstinline!r + 1!}；当前模板已在末段结束时退出。

\hypertarget{ux675cux6559ux7b5b}{%
\paragraph{杜教筛}\label{ux675cux6559ux7b5b}}

记

\[
M(n)=\sum_{i=1}^{n}\mu(i),\qquad
\Phi(n)=\sum_{i=1}^{n}\varphi(i).
\]

则对 \(n\ge1\)，有

\[
M(n)=1-\sum_{l=2}^{n}M\left(\left\lfloor\frac{n}{l}\right\rfloor\right),
\]

\[
\Phi(n)=\frac{n(n+1)}{2}
-\sum_{l=2}^{n}\Phi\left(\left\lfloor\frac{n}{l}\right\rfloor\right).
\]

右侧可以按商分块，同一段的函数值只计算一次，再乘段长。小范围预筛，大范围记忆化。预处理上限 \(B\) 应根据查询规模和内存选择；取 \(B\) 约为 \(n^{2/3}\) 时，常用的总时间界为 \(O(n^{2/3})\)。固定一个很小的 \(B\) 后，不能继续对任意大的查询使用这个时间估计。

\hypertarget{ux6700ux5927ux516cux7ea6ux6570ux5377ux79efux4e0eux6700ux5c0fux516cux500dux6570ux5377ux79ef}{%
\paragraph{最大公约数卷积与最小公倍数卷积}\label{ux6700ux5927ux516cux7ea6ux6570ux5377ux79efux4e0eux6700ux5c0fux516cux500dux6570ux5377ux79ef}}

倍数 zeta 变换为

\[
A_d=\sum_{d\mid i}a_i.
\]

对两个数组分别变换、逐项相乘，再做逆变换，可以得到

\[
c_k=\sum_{\gcd(i,j)=k}a_i b_j.
\]

约数 zeta 变换为

\[
A_x=\sum_{i\mid x}a_i.
\]

同样先变换、相乘，再做逆变换，可以得到

\[
c_k=\sum_{\operatorname{lcm}(i,j)=k}a_i b_j.
\]

当前实现直接枚举整除关系，时间为 \(O(N\log N)\)，不依赖额外筛表。数组下标从 \(1\) 开始，\passthrough{\lstinline!a[0]!} 不参与。最小公倍数卷积只返回下标不超过 \(N\) 的部分。

原地变换时，循环方向决定某个值是否已经被累加过，不能随意交换。

\hypertarget{ux7d20ux6570ux5e42ux7ec4ux5408ux6570ux4e0eux6a21ux5e73ux65b9ux6839}{%
\subsubsection{5. 素数幂组合数与模平方根}\label{ux7d20ux6570ux5e42ux7ec4ux5408ux6570ux4e0eux6a21ux5e73ux65b9ux6839}}

\hypertarget{exlucas}{%
\paragraph{exLucas}\label{exlucas}}

将模数分解成互质的素数幂 \(p^e\)，分别计算组合数，再用中国剩余定理合并。每个素数幂下的步骤是：

\begin{enumerate}
\def\labelenumi{\arabic{enumi}.}
\tightlist
\item
  计算 \(n!\)、\(r!\) 和 \((n-r)!\) 中素因子 \(p\) 的指数差。
\item
  若指数差不小于 \(e\)，结果模 \(p^e\) 为 \(0\)。
\item
  从阶乘中去掉所有 \(p\) 因子，对剩余部分求逆，再乘回所需的 \(p\) 的幂。
\end{enumerate}

不能对含 \(p\) 因子的普通阶乘直接求逆。当前实现为每个素数幂预处理完整表，因此时间和空间为 \(O\!\left(\sum p^e\right)\)。如果模数本身就是一个很大的素数，表也会很大，并不会因为只有一个因子而节省空间。

\hypertarget{ux6a21ux5e73ux65b9ux6839}{%
\paragraph{模平方根}\label{ux6a21ux5e73ux65b9ux6839}}

模平方根求解

\[
x^2\equiv a\quad(\operatorname{mod} p).
\]

当前 Tonelli-Shanks 模板要求 \(p\) 为质数。先处理 \(p=2\) 和 \(a=0\)，再用 Euler 判别法判断是否有解。非零根成对出现，为 \(x\) 和 \(p-x\)。接口返回较小的非负根，无解返回 \passthrough{\lstinline!-1!}。

Miller-Rabin 和 Pollard-Rho 的接口使用有符号 \passthrough{\lstinline!ll!}，没有覆盖 \passthrough{\lstinline!u64!} 的全部范围。Pollard-Rho 是随机算法，对每个输入没有统一的确定运行时间保证。

\hypertarget{ux7ebfux6027ux4ee3ux6570ux4e0eux5f02ux6216ux7ebfux6027ux57fa}{%
\subsubsection{6. 线性代数与异或线性基}\label{ux7ebfux6027ux4ee3ux6570ux4e0eux5f02ux6216ux7ebfux6027ux57fa}}

\hypertarget{ux5f02ux6216ux9ad8ux65afux6d88ux5143}{%
\paragraph{异或高斯消元}\label{ux5f02ux6216ux9ad8ux65afux6d88ux5143}}

在 \(\mathbb F_2\) 上，加减法都是异或。若方程组有 \(m\) 个未知数、秩为 \(r\)，且有解，则解的数量为 \(2^{m-r}\)。将自由变量设为 \(0\)，可以得到一个特解。

当前接口把 \passthrough{\lstinline!bitset!} 的第 \passthrough{\lstinline!m!} 位用作常数项；变量位是从第 \passthrough{\lstinline!0!} 位到第 \passthrough{\lstinline!m - 1!} 位。

\hypertarget{ux5f02ux6216ux7ebfux6027ux57fa}{%
\paragraph{异或线性基}\label{ux5f02ux6216ux7ebfux6027ux57fa}}

异或线性基表示所有子集异或值构成的向量空间。设原集合有 \(n\) 个元素，线性基的秩为 \(r\)：

\begin{itemize}
\tightlist
\item
  不同异或值有 \(2^r\) 个，包含空集得到的 \(0\)。
\item
  每个可以表示的值都有 \(2^{n-r}\) 种子集表示，这里包含空集。
\item
  插入过程中出现线性相关，当且仅当存在非空子集的异或为 \(0\)。
\end{itemize}

第 \(k\) 小异或值按不同数值排序，不按子集表示的重复次数排序。当前 \passthrough{\lstinline!kth!} 从 \passthrough{\lstinline!0!} 开始编号；设置 \passthrough{\lstinline!nonempty=true!} 后排除空集。若仍有非空子集的异或为 \(0\)，结果中依然保留 \(0\)。

完整的 \(64\) 位线性基可能有 \(2^{64}\) 个不同值，不能用 \passthrough{\lstinline!u64!} 存放这个总数。实现中也不能计算 \passthrough{\lstinline!1ULL << 64!}。

\hypertarget{ux666eux901aux6d88ux5143ux7684ux9002ux7528ux8303ux56f4}{%
\paragraph{普通消元的适用范围}\label{ux666eux901aux6d88ux5143ux7684ux9002ux7528ux8303ux56f4}}

当前 \passthrough{\lstinline!matrix\_inv!} 和 \passthrough{\lstinline!determinant!} 使用域上的精确消元，主要用于质数模。不能把其中的除法直接换成普通整数除法，也不能无条件用于合数模。浮点版本则需要另外处理主元选择和误差判断。

\hypertarget{ux7ebfux6027ux9012ux63a8ux751fux6210ux51fdux6570ux4e0eux63d2ux503c}{%
\subsubsection{7. 线性递推、生成函数与插值}\label{ux7ebfux6027ux9012ux63a8ux751fux6210ux51fdux6570ux4e0eux63d2ux503c}}

\hypertarget{ux666eux901aux751fux6210ux51fdux6570ux4e0eux6307ux6570ux751fux6210ux51fdux6570}{%
\paragraph{普通生成函数与指数生成函数}\label{ux666eux901aux751fux6210ux51fdux6570ux4e0eux6307ux6570ux751fux6210ux51fdux6570}}

记 \([x^n]F(x)\) 为 \(F(x)\) 的 \(x^n\) 项系数。普通生成函数使用 \(A(x)=\sum_{n\ge0}a_nx^n\)，于是

\[
[x^n]A(x)B(x)=\sum_{k=0}^{n}a_kb_{n-k},\qquad
[x^n]\frac{A(x)}{1-x}=\sum_{k=0}^{n}a_k.
\]

若 \(B(x)\) 表示一段非空结构，且 \(B(0)=0\)，则由任意多个有顺序的段组成的结构，其普通生成函数为 \(1/(1-B(x))\)。

指数生成函数使用 \(A(x)=\sum_{n\ge0}a_nx^n/n!\)。若 \(A,B\) 都按此形式定义，则

\[
n![x^n]A(x)B(x)=\sum_{k=0}^{n}\binom{n}{k}a_kb_{n-k}.
\]

这里的组合数表示先选择哪些标号分给第一部分。若 \(B(x)\) 统计非空的标号连通块，则由这些块组成的无序集合使用 \(\exp(B(x))\)，其中 \(B(0)=0\)；块的计数和对称性必须已经正确计入 \(B(x)\)。

这些运算通常视为形式幂级数，不要求给 \(x\) 代入实数，也不依赖解析收敛。在模意义下使用阶乘、积分或指数时，还需保证涉及的整数分母可逆。

\hypertarget{lagrange-ux53cdux6f14}{%
\paragraph{Lagrange 反演}\label{lagrange-ux53cdux6f14}}

设 \(F(0)=0\)，且 \(F(x)=x\phi(F(x))\)。对于 \(n\ge k\ge1\)，有

\[
[x^n]F(x)^k=\frac{k}{n}[t^{n-k}]\phi(t)^n.
\]

它适合处理``一个根连接若干子结构''产生的隐式生成函数。若 \(T=x\exp(T)\)，则 \([x^n]T=n^{n-1}/n!\)，对应 \(n\) 个标号点的有根树。

该形式在有理数或特征零的域上使用最直接。模意义下不能在 \(n\) 不可逆时直接除以 \(n\)；应先在整数意义下变形，或使用适用的实现。

\hypertarget{ux9012ux63a8ux7cfbux6570ux7684ux987aux5e8f}{%
\paragraph{递推系数的顺序}\label{ux9012ux63a8ux7cfbux6570ux7684ux987aux5e8f}}

所有递推接口都使用

\[
a_i=\sum_{j=0}^{k-1}c_j a_{i-j-1}.
\]

即 \passthrough{\lstinline!c[0]!} 对应最近一项，项编号从 \passthrough{\lstinline!0!} 开始。

\hypertarget{kitamasabm-ux4e0e-bostan-mori}{%
\paragraph{Kitamasa、BM 与 Bostan-Mori}\label{kitamasabm-ux4e0e-bostan-mori}}

Kitamasa 将 \(x^n\) 对特征多项式

\[
x^k-c_0x^{k-1}-c_1x^{k-2}-\cdots-c_{k-1}
\]

取模，从而把第 \(n\) 项写成前 \(k\) 项的线性组合。时间为 \(O(k^2\log n)\)，只需要加减乘，因此可以用于合数模。

BM 在域上从已知前缀寻找最短线性递推。若已经知道真实阶数不超过 \(k\)，前 \(2k\) 项足以恢复递推。仅凭一段任意数列猜出递推，并不能证明它对后续成立；应使用额外项检查，并结合题目结构证明阶数上界。

Bostan-Mori 从生成函数出发。设

\[
A(x)=\sum_{i\ge0}a_i x^i,\qquad
Q(x)=1-c_0x-\cdots-c_{k-1}x^k.
\]

令 \(P(x)\) 为 \(A(x)Q(x)\) 的前 \(k\) 项，就有 \(A(x)=P(x)/Q(x)\)。分子分母同乘 \(Q(-x)\) 后，分母只剩偶数次项。根据 \(n\) 的奇偶保留分子的对应项，再将下标减半。使用 NTT 后，时间为 \(O(k\log k\log n)\)。

一个数列在不同模数下的最短递推阶数可能不同。不能把单个模数下 BM 得到的结果直接当成其他模数或整数范围内的结论。

\hypertarget{ux8fdeux7eedux70b9ux63d2ux503cux4e0eux5e42ux548c}{%
\paragraph{连续点插值与幂和}\label{ux8fdeux7eedux70b9ux63d2ux503cux4e0eux5e42ux548c}}

已知 \(f(0),\ldots,f(k)\)，且 \(\deg f\le k\)，可用前后缀积在 \(O(k)\) 时间内求出一个新点的值。要求这些横坐标在域中互异；在质数模 \(P\) 下，\(k<P\) 是充分条件。

幂和

\[
\sum_{i=1}^{n}i^k
\]

可以表示为关于 \(n\) 的 \(k+1\) 次多项式。当前 \passthrough{\lstinline!power\_sum!} 使用横坐标从 \(0\) 到 \(k+1\) 的样本，要求 \(k+1<P\)。生成样本还需要 \(O(k\log(k+1))\) 时间。

\hypertarget{taylor-shift}{%
\paragraph{Taylor shift}\label{taylor-shift}}

给定 \(f(x)=\sum_{j\ge0}a_jx^j\)，需要计算 \(f(x+c)\)。其第 \(i\) 项系数为

\[
\sum_{j\ge i}a_j\binom{j}{i}c^{j-i}.
\]

乘上阶乘并反转一侧后，可转化成一次卷积。使用时需满足阶乘可逆和 NTT 长度条件。

\hypertarget{ux5b50ux96c6ux5377ux79efux4e0e-or-ux5377ux79ef}{%
\paragraph{子集卷积与 OR 卷积}\label{ux5b50ux96c6ux5377ux79efux4e0e-or-ux5377ux79ef}}

OR 卷积允许两个集合重叠；子集卷积只合并不相交集合：

\[
c_S=\sum_{T\subseteq S}a_T b_{S\setminus T}.
\]

按集合大小分层，分别做子集 zeta 变换，再进行层间卷积和逆变换，最后只取层数等于集合大小的结果。时间为 \(O(k^2 2^k)\)，空间为 \(O(k2^k)\)。

当 \(k=20\) 且元素类型为 \passthrough{\lstinline!mint!} 时，当前实现的主要数组约占 \(168\,\mathrm{MiB}\)，还没有计入少量额外数组。

\hypertarget{nimsg-ux4e0eux53cdux5e38ux6e38ux620f}{%
\subsubsection{8. Nim、SG 与反常游戏}\label{nimsg-ux4e0eux53cdux5e38ux6e38ux620f}}

普通 Nim 每次从一堆中取走任意正数颗石子，不能行动者输。所有堆大小的异或和非零，当且仅当当前局面必胜。

若异或和为 \(s\ne0\)，可选择在 \(s\) 的最高位上为 \(1\) 的一堆，把它的大小 \(x\) 改为 \(x\mathbin{\oplus}s\)。这个数严格小于 \(x\)，且操作后总异或和为 \(0\)。

对于有限、无环的公平组合游戏，定义

\[
\operatorname{sg}(u)=\operatorname{mex}
\{\operatorname{sg}(v):u\to v\}.
\]

其中 \(\operatorname{mex}\) 是集合中没有出现的最小非负整数。多个独立子游戏的 SG 值按位异或，就得到它们组合后的 SG 值。

这里要求双方在同一局面下的可选动作相同，并采用不能行动者输的规则。有环图可能产生和局，不能直接套用当前 DAG 模板。

反常 Nim 采用拿走最后一颗者输的规则：

\begin{itemize}
\tightlist
\item
  若存在大于 \(1\) 的堆，仍由异或和是否非零判断胜负。
\item
  若所有堆的大小都是 \(0\) 或 \(1\)，非空堆数为偶数则胜，为奇数则负。
\item
  当前接口把没有石子的局面视为当前玩家胜，这是反常规则下的递归终局约定。
\end{itemize}

反常 Nim 的结论不能推广为把一般 SG 游戏的胜负直接反转。

\hypertarget{bash-ux535aux5f08ux4e0e-wythoff-ux535aux5f08}{%
\paragraph{Bash 博弈与 Wythoff 博弈}\label{bash-ux535aux5f08ux4e0e-wythoff-ux535aux5f08}}

Bash 博弈只有一堆 \(n\) 个石子，每次取 \(1\) 到 \(m\) 个，其中 \(m\ge1\)，不能行动者输。必败态为 \(n\equiv0\pmod{m+1}\)。必胜时先取 \(n\bmod(m+1)\) 个，之后与对手每轮合取 \(m+1\) 个。

Wythoff 博弈有两堆石子，每次可以从任意一堆取正数个，或从两堆各取相同的正数个，不能行动者输。将堆大小排序为 \(0\le a\le b\)，令 \(k=b-a\)、\(\phi=(1+\sqrt5)/2\)，必败态恰为

\[
(a,b)=\left(\lfloor k\phi\rfloor,\lfloor k\phi\rfloor+k\right),\qquad k\ge0.
\]

大整数范围下，浮点计算 \(\lfloor k\phi\rfloor\) 可能差一。可用精确整数平方根计算

\[
\lfloor k\phi\rfloor
=\left\lfloor\frac{k+\lfloor\sqrt{5k^2}\rfloor}{2}\right\rfloor,
\]

但必须使用能容纳 \(5k^2\) 的整数类型。这是必败态判定公式，不能将它直接当作一般局面的 SG 值。

\hypertarget{ux6570ux503cux79efux5206}{%
\subsubsection{9. 数值积分}\label{ux6570ux503cux79efux5206}}

Simpson 公式对三次及以下多项式精确。设一整个区间的估计为 \(S\)，分成两段后的估计为 \(S_1,S_2\)，则用

\[
\Delta=S_1+S_2-S
\]

判断是否继续细分，并用

\[
S_1+S_2+\frac{\Delta}{15}
\]

作为修正后的估计。

\passthrough{\lstinline!eps!} 表示目标绝对误差，递归时将误差额度平分给两段。这个估计不能保证对任意函数都准确：高频振荡可能被采样点漏掉；已知函数变化尺度时，应主动分段。

不连续点应分段处理，端点存在无穷大时应换元或使用专门方法。达到递归深度上限，或浮点区间已无法继续细分时，模板返回当前估计，并不表示已经满足误差目标。

积分方向反转时，答案取负；零长区间直接返回 \(0\)。

\hypertarget{ux6982ux7387ux4e0eux671fux671b}{%
\subsubsection{10. 概率与期望}\label{ux6982ux7387ux4e0eux671fux671b}}

\hypertarget{ux6307ux793aux53d8ux91cfux5c3eux548cux516cux5f0fux4e0eux65b9ux5dee}{%
\paragraph{指示变量、尾和公式与方差}\label{ux6307ux793aux53d8ux91cfux5c3eux548cux516cux5f0fux4e0eux65b9ux5dee}}

令 \(I_A\) 表示事件 \(A\) 是否发生，则 \(\mathbb E[I_A]=\Pr(A)\)。期望满足线性性：

\[
\mathbb E\!\left[\sum_i X_i\right]=\sum_i\mathbb E[X_i].
\]

有限求和不要求变量独立；非负项的无限求和也可以交换求和与期望。计算``出现多少次''``有多少对满足条件''时，通常先拆成指示变量。

非负整数随机变量 \(X\) 满足

\[
\mathbb E[X]=\sum_{k\ge1}\Pr(X\ge k),\qquad
\mathbb E[X^2]=\sum_{k\ge1}(2k-1)\Pr(X\ge k).
\]

方差满足

\[
\operatorname{Var}(X)=\mathbb E[X^2]-\mathbb E[X]^2.
\]

独立变量的方差可以相加；不独立时需要计入协方差。不能由期望的线性性推出 \(\mathbb E[XY]=\mathbb E[X]\mathbb E[Y]\)，也不能推出 \(\mathbb E[1/X]=1/\mathbb E[X]\)。

\hypertarget{ux51e0ux4f55ux5206ux5e03ux96c6ux9f50ux95eeux9898ux4e0eux968fux673aux6392ux5217}{%
\paragraph{几何分布、集齐问题与随机排列}\label{ux51e0ux4f55ux5206ux5e03ux96c6ux9f50ux95eeux9898ux4e0eux968fux673aux6392ux5217}}

每次独立尝试的成功概率均为 \(p>0\)，计入首次成功的那次尝试，期望尝试次数为 \(1/p\)。若只统计成功前的失败次数，期望为 \((1-p)/p\)。

独立、等概率地抽取 \(n\ge1\) 种物品，集齐全部种类的期望次数为

\[
nH_n,\qquad H_n=\sum_{k=1}^{n}\frac1k.
\]

因为已有 \(i\) 种时，获得新品种的概率为 \((n-i)/n\)。若各类出现概率不同，不能继续套用 \(nH_n\)。

均匀随机的 \(n\) 元排列中，固定点数的期望为 \(1\)（\(n\ge1\)），逆序对数的期望为 \(n(n-1)/4\)，置换环数的期望为 \(H_n\)（\(n\ge1\)）。分别考虑每个位置、每对位置，以及按编号插入元素时是否新开一个环即可推导。

\hypertarget{ux671fux671b-dp-ux4e0eux6709ux73afux72b6ux6001}{%
\paragraph{期望 DP 与有环状态}\label{ux671fux671b-dp-ux4e0eux6709ux73afux72b6ux6001}}

若从状态 \(u\) 走向 \(v\) 的概率为 \(p_{u,v}\)，本次转移的代价为 \(c_{u,v}\)，则到达终态的期望总代价满足

\[
E_u=\sum_v p_{u,v}(c_{u,v}+E_v),\qquad E_{\mathrm{terminal}}=0.
\]

无环状态可逆拓扑计算。有环时通常要解线性方程组；只有自环时，可先把含 \(E_u\) 的项移到左侧，再除以 \(1-p_{u,u}\)。若无法保证以概率一到达终态，或期望不有限，不能将方程组的某个解直接当作期望。

模意义下的概率表示要求相关分母可逆，且并不保留大小关系。浮点方程则需要处理主元选择与误差。

\hypertarget{ux62bdux5c49ux539fux7406ux4e0eux504fux5e8f}{%
\subsubsection{11. 抽屉原理与偏序}\label{ux62bdux5c49ux539fux7406ux4e0eux504fux5e8f}}

\hypertarget{ux62bdux5c49ux539fux7406ux4e0eux524dux7f00ux548c}{%
\paragraph{抽屉原理与前缀和}\label{ux62bdux5c49ux539fux7406ux4e0eux524dux7f00ux548c}}

把 \(n\) 个物体放入 \(m\ge1\) 个抽屉，至少有一个抽屉含不少于 \(\lceil n/m\rceil\) 个物体。

任意 \(m\) 个整数中，都存在一个非空连续子段，其和能被 \(m\) 整除：考虑包含空前缀的 \(m+1\) 个前缀和，它们模 \(m\) 只有 \(m\) 种余数，必有两个相同。

\hypertarget{dilworth-ux4e0e-mirsky-ux5b9aux7406}{%
\paragraph{Dilworth 与 Mirsky 定理}\label{dilworth-ux4e0e-mirsky-ux5b9aux7406}}

在有限偏序集中，链是任意两个元素都可比较的子集，反链是任意两个不同元素都不可比较的子集。两个定理分别给出

\[
\text{最大反链大小}=\text{划分成链的最少条数},
\]

\[
\text{最长链长度}=\text{划分成反链的最少份数}.
\]

前者是 Dilworth 定理，后者是 Mirsky 定理。这里都要求各份覆盖全部元素，且每个元素恰好属于一份。

求最少链划分时，可将每个元素拆成左右两个点，对每一对严格偏序关系 \(u<v\) 连边，答案为 \(n-|M|\)，其中 \(M\) 是最大二分图匹配。需要使用全部可比关系，不能只取一个给定 DAG 的原始边就认为已经建立了偏序关系。

对于数列，可按``下标较小且数值不大于''定义偏序：最少划分成多少个不下降子序列，等于最长严格下降子序列的长度。若改用严格上升子序列，另一侧就对应最长不上升子序列；重复值决定这里能否取等号。

\hypertarget{ux53c2ux8003ux8d44ux6599}{%
\subsubsection{12. 参考资料}\label{ux53c2ux8003ux8d44ux6599}}

\hypertarget{ux6570ux5b66ux7ed3ux8bba}{%
\paragraph{数学结论}\label{ux6570ux5b66ux7ed3ux8bba}}

\begin{itemize}
\tightlist
\item
  \href{https://dlmf.nist.gov/27.14}{DLMF：整数拆分与五边形数定理}：生成函数、广义五边形数与拆分递推。
\item
  \href{https://dlmf.nist.gov/26.3}{DLMF：二项式系数}：组合恒等式。
\item
  \href{https://aofa.cs.princeton.edu/30gf/}{Princeton：生成函数}：普通生成函数、指数生成函数与 Lagrange 反演。
\item
  \href{https://crypto.stanford.edu/pbc/notes/numbertheory/order.html}{Stanford：模乘法群中元素的阶}：阶与 Euler 定理。
\item
  \href{https://ocw.mit.edu/courses/18-600-probability-and-random-variables-fall-2019/pages/lecture-notes/}{MIT：概率与随机变量讲义}：期望、方差与常见概率模型。
\item
  \href{https://mathcircle.berkeley.edu/sites/default/files/handouts/2019/BMC-mashup-debruijn-order-dilworth-hall.pdf}{Berkeley Math Circle：偏序、Dilworth 与 Mirsky 定理}。
\item
  \href{https://arxiv.org/abs/2310.07735}{Wythoff 博弈与黄金比例的推导}。
\end{itemize}

\hypertarget{ux6a21ux677fux5b9eux73b0}{%
\paragraph{模板实现}\label{ux6a21ux677fux5b9eux73b0}}

实现时对比了以下资料，并按本库的命名和已有接口整理：

\begin{itemize}
\tightlist
\item
  \href{https://github.com/kth-competitive-programming/kactl/blob/main/content/numerical/BerlekampMassey.h}{KACTL 的 BM}、\href{https://github.com/kth-competitive-programming/kactl/blob/main/content/numerical/LinearRecurrence.h}{线性递推}和\href{https://github.com/kth-competitive-programming/kactl/blob/main/content/numerical/IntegrateAdaptive.h}{自适应积分}：参考独立短函数的组织方式；积分另外缓存函数值并设置递归深度上限。
\item
  \href{https://github.com/NyaanNyaan/library/blob/master/fps/berlekamp-massey.hpp}{Nyaan 的 BM}：对比递推系数的返回形式，本库统一将最近一项的系数放在前面。
\item
  \href{https://github.com/ei1333/library/blob/master/math/fps/formal-power-series-friendly-ntt.hpp}{ei1333 的 FPS}：参考多项式算法，仍使用本库已有的 \passthrough{\lstinline!convolution!} 和独立函数接口。
\item
  \href{https://github.com/atcoder/ac-library/blob/master/atcoder/math.hpp}{ACL 的数学接口}和 \href{https://github.com/atcoder/ac-library/blob/master/atcoder/modint.hpp}{ModInt}：参考参数条件的表达方式，继续使用本库已有的 \passthrough{\lstinline!ModInt!}。
\end{itemize}
